\documentclass[11pt,a4paper,oneside]{report}

% --- Korean & Fonts ---
\usepackage[hangul]{kotex}
\setmainhangulfont{Malgun Gothic}
\setsanshangulfont{Malgun Gothic}
\setmonohangulfont{Malgun Gothic}

% --- Math ---
\usepackage{amsmath, amssymb, amsthm, mathtools}
\usepackage{bm}

% --- Layout ---
\usepackage[a4paper,margin=2.3cm,top=2.5cm,bottom=2.5cm]{geometry}
\usepackage{setspace}
\linespread{1.15}

% --- Tables, figures ---
\usepackage{booktabs, array, tabularx, longtable}
\usepackage{graphicx}
\usepackage{caption}
\captionsetup{font=small}

% --- Code listings ---
\usepackage{listings}
\usepackage{xcolor}
\definecolor{codebg}{RGB}{248,248,245}
\definecolor{codecomment}{RGB}{106,153,85}
\definecolor{codekw}{RGB}{0,0,180}
\definecolor{codestr}{RGB}{163,21,21}
\definecolor{codenum}{RGB}{120,120,120}
\definecolor{coderule}{RGB}{210,210,210}
\lstdefinestyle{cpp}{
  language=C++,
  basicstyle=\ttfamily\footnotesize,
  keywordstyle=\color{codekw}\bfseries,
  commentstyle=\color{codecomment}\itshape,
  stringstyle=\color{codestr},
  numbers=left,
  numberstyle=\tiny\color{codenum},
  backgroundcolor=\color{codebg},
  frame=single,
  rulecolor=\color{coderule},
  showstringspaces=false,
  breaklines=true,
  tabsize=2,
  captionpos=b,
  morekeywords={Eigen,Matrix3d,Vector3d,Matrix4d,Quaterniond,VectorXd,MatrixXd,Mat15,Vec15,size_t,uchar,ImuData,CamData,GtData,CameraParams,ImuNoiseParams,ImuPropagator,LcEkf,StereoTracker,Pose,KittiRawReader,YAML,Node,fs,path,cv}
}
\lstdefinestyle{shell}{
  basicstyle=\ttfamily\footnotesize,
  backgroundcolor=\color{codebg},
  frame=single,
  rulecolor=\color{coderule},
  showstringspaces=false,
  breaklines=true,
  tabsize=2,
  captionpos=b
}
\lstdefinestyle{yaml}{
  basicstyle=\ttfamily\footnotesize,
  backgroundcolor=\color{codebg},
  frame=single,
  rulecolor=\color{coderule},
  commentstyle=\color{codecomment}\itshape,
  morecomment=[l]{\#},
  showstringspaces=false,
  breaklines=true,
  captionpos=b
}
\lstdefinestyle{py}{
  language=Python,
  basicstyle=\ttfamily\footnotesize,
  keywordstyle=\color{codekw}\bfseries,
  commentstyle=\color{codecomment}\itshape,
  stringstyle=\color{codestr},
  backgroundcolor=\color{codebg},
  frame=single,
  rulecolor=\color{coderule},
  showstringspaces=false,
  breaklines=true,
  tabsize=2,
  captionpos=b
}
\lstset{style=cpp}

% --- Hyperref ---
\usepackage[colorlinks=true,linkcolor=blue!50!black,urlcolor=blue!60!black,citecolor=blue!50!black]{hyperref}

% --- TikZ ---
\usepackage{tikz}
\usetikzlibrary{positioning,arrows.meta,shapes.geometric,fit,backgrounds,calc}

% --- Boxes ---
\usepackage[most]{tcolorbox}
\newtcolorbox{kbox}[1][]{colback=blue!4,colframe=blue!40!black,sharp corners,boxrule=0.5pt,#1}
\newtcolorbox{wbox}[1][]{colback=red!4,colframe=red!40!black,sharp corners,boxrule=0.5pt,#1}
\newtcolorbox{notebox}[1][]{colback=yellow!8,colframe=orange!60!black,sharp corners,boxrule=0.5pt,#1}

% --- Math shortcuts ---
\newcommand{\R}{\mathbb{R}}
\newcommand{\so}{\mathfrak{so}}
\newcommand{\SO}{\mathrm{SO}}
\newcommand{\SE}{\mathrm{SE}}
\newcommand{\T}{^{\!\top}}
\newcommand{\skewop}[1]{[#1]_{\!\times}}
\DeclareMathOperator{\Exp}{Exp}
\DeclareMathOperator{\Log}{Log}
\DeclareMathOperator{\diag}{diag}

% --- Title ---
\title{\textbf{LC-EKF VIO on KITTI raw}\\[0.3em]
\large 백지에서 시작하는 풀 파이프라인 가이드 \\
\large {\small 개념 · 수학 유도 · 코드 분석 · 운용 절차}}
\author{C++ ROS-free Stereo VIO 학습 노트}
\date{2026}

\begin{document}

\maketitle

\begin{abstract}
\noindent
이 문서는 \texttt{06\_cpp\_LC-EKF\_VIO\_kitti\_raw} 프로젝트를 \textbf{외부 도움 없이} 처음부터 끝까지 굴려 보고자 하는 사용자를 위한 통합 가이드다. 다음 네 가지 층위를 한 권에 묶었다:

\begin{enumerate}
\item \textbf{개념 (Why)}: VIO/SLAM/EKF 가 무엇이고 왜 필요한가, 왜 스테레오 + IMU 조합인가, 왜 LC-EKF 구조인가.
\item \textbf{수학 (What)}: SO(3) 회전 표현, error-state EKF, RK4 적분, F/G/Q 행렬의 유도, Kalman gain 의 Joseph form, Stereo triangulation, PnP, Umeyama 정합. 각 수식은 \emph{왜 그렇게 쓰는가} 까지 짚는다.
\item \textbf{코드 (How)}: 3개 코어 모듈(\texttt{StereoTracker}, \texttt{ImuPropagator}, \texttt{LcEkf})과 2개 운용 모듈(\texttt{KittiRawReader}, \texttt{run\_vio.cpp})에 대한 줄 단위 분석과 수식-코드 1:1 매핑.
\item \textbf{운용 (Do)}: KITTI 데이터셋 다운로드, 캘리브레이션 추출, YAML 설정, 빌드, 실행, 시각화, 다중 시퀀스 일괄 테스트, 트러블슈팅.
\end{enumerate}

본 가이드를 처음부터 끝까지 따라가면 약 1시간 30분 내에 첫 trajectory plot 을 얻을 수 있고, 그 후 후반부의 튜닝 시나리오로 LC-EKF 의 한계와 가능성을 직접 실험해 볼 수 있다.
\end{abstract}

\tableofcontents
\clearpage

% =====================================================================
\chapter{30분 개관 — 백지에서 출발하는 이해}
% =====================================================================

이 장은 VIO·SLAM·EKF·KITTI 를 \textbf{처음 들어보는 사람}이 백지 상태에서 ``이게 뭐 하는 거고 내가 뭘 할지'' 를 30분 안에 잡도록 쓴 도입부다.

\section{왜 이런 시스템이 필요한가 — 위치를 안다는 문제}

자율주행/드론/로봇청소기 모두 ``내가 지금 어디 있는가''를 매 0.1 초마다 알아야 한다. 방법별 정밀도/한계:

\begin{center}\small
\begin{tabular}{lll}
\toprule
방법 & 정밀도 & 한계 \\
\midrule
GPS (스마트폰) & 1 \textasciitilde{} 10 m & 터널/실내/숲 무력. 1 Hz 라 차량 동역학에 너무 느림 \\
GPS-RTK (정밀) & 1 \textasciitilde{} 10 cm & 기지국 인프라 필요. GPS 신호 끊기면 똑같이 무력 \\
휠 오도메트리 & 단기 정밀 & 미끄러짐/타이어 압력으로 누적 오차 \\
\textbf{VIO (이 프로젝트)} & 누적 \textbf{0.5\textasciitilde{}5\%} & GPS 없이도 동작, 인프라 불필요, 실시간 \\
\bottomrule
\end{tabular}
\end{center}

\textbf{VIO = Visual-Inertial Odometry} — 카메라(visual)와 IMU(inertial) 만으로 \textbf{6-DoF pose}(3D 위치 + 3D 회전)를 매 순간 추정한다. GPS 의존을 줄이고 실내에서도 동작하며, 카메라 + 센서만 있으면 어디서든 자세 추적이 가능하다.

\section{왜 카메라랑 IMU를 같이 쓰는가}

각 센서 단독의 약점이 명확하고 서로 \textbf{상보}이기 때문이다.

\begin{center}\small
\begin{tabular}{p{0.32\textwidth}p{0.32\textwidth}p{0.32\textwidth}}
\toprule
& \textbf{카메라(스테레오) 단독} & \textbf{IMU 단독} \\
\midrule
장점 &
정밀한 절대 위치 (특징점 충분 시 0.1 m). 미터 스케일 자동. &
100\textasciitilde{}200 Hz 고주파. 빛/날씨/텍스처 무관. \\
\midrule
단점 &
빠른 회전 시 모션 블러로 추적 실패. 어두운 터널, 단조 텍스처에서 약함. 10 Hz 한계. &
\textbf{누적 오차로 발산} (가속도 두 번 적분 → 3\textasciitilde{}5초에 1 m, 1분에 수십\textasciitilde{}수백 m). \\
\bottomrule
\end{tabular}
\end{center}

\textbf{둘을 합치면}: 카메라가 절대 기준 역할을 해 IMU 드리프트를 매 프레임 잡아주고, IMU 가 카메라 사이의 빈 시간을 부드럽게 채워 준다. 한쪽이 일시 실패해도 다른 쪽이 버틴다.

\section{왜 스테레오인가 — 미터 스케일의 출처}

모노 카메라 한 개의 본질적 한계: 50 cm 떨어진 작은 사과와 5 m 떨어진 큰 사과를 픽셀로는 구별할 수 없다. 이를 \textbf{스케일 모호성(scale ambiguity)} 이라 한다.

스테레오 두 카메라:
\begin{itemize}
\item 좌\(\leftrightarrow\)우 픽셀 차이(\emph{disparity}) + 미리 측정한 두 카메라 사이 거리(\textbf{baseline})
\item 깊이 공식: \(\boxed{Z = \dfrac{f_x \cdot b}{d}}\)
\item 우리 KITTI calib 의 baseline = \textbf{53.7 cm} (KITTI 측이 레이저로 잰 물리값)
\end{itemize}

이 한 값에서 모든 미터 스케일이 파생된다. 좌측이 \texttt{cam0} (\texttt{image\_00}), 우측이 \texttt{cam1} (\texttt{image\_01}). 컬러 페어(\texttt{image\_02/03})도 있지만 본 코드는 \textbf{흑백 스테레오 페어} 만 사용한다 — KLT 추적이 컬러 보간 잡음에 덜 흔들리기 때문이다.

\section{왜 EKF인가 — 두 측정값을 통계적으로 섞기}

\textbf{EKF (Extended Kalman Filter)} 는 비선형 시스템에서 각 센서의 잡음 모델을 들고 가중평균하는 표준 도구다.

\begin{kbox}
\textit{IMU 적분으로 위치는 \((10.0, 5.0, 0.1)\) 일 것이다 (불확실성 \(\pm 0.5\) m)} \\
\textit{VO가 측정한 위치는 \((10.5, 4.7, 0.2)\) 다 (잡음 \(\pm 0.1\) m)} \\
\(\downarrow\) EKF 가 \(\sigma\) 로 가중해 합침 \\
\textit{통합 추정: \((10.45, 4.73, 0.19)\)} \quad ← VO 에 더 가중치 (잡음 작음)
\end{kbox}

내부적으로 \textbf{15차원 상태} (위치 3 + 속도 3 + 자세 3 + 자이로 bias 3 + 가속도 bias 3) 와 그에 대응하는 \(15 \times 15\) 공분산 행렬 \(\mathbf{P}\) 를 들고 다닌다.

본 프로젝트는 \textbf{LC-EKF (Loosely-Coupled)} 구조 — VO 의 \emph{출력 위치}를 측정값으로 받는다. 더 정밀한 변종은 \textbf{Tight MSCKF} (특징점 픽셀을 직접 측정값으로) 로, 본 프로젝트의 다음 마일스톤이다.

\section{KITTI 데이터셋이 뭐고 왜 쓰나}

\begin{itemize}
\item 독일 KIT (카를스루에 공과대) + 미국 TTIC (토요타 시카고 연구소) 가 2012년 공개
\item 차량에 장착: 흑백 스테레오 + 컬러 스테레오 + Velodyne LiDAR + \textbf{OXTS RT3003 GPS/IMU 통합 유닛}
\item 도시/주거지/도로/고속도로/캠퍼스 다양한 환경 50 시간 +
\item \textbf{VIO/SLAM 알고리즘 평가의 사실상 표준 벤치마크}
\end{itemize}

본 프로젝트가 쓰는 부분:
\begin{itemize}
\item \texttt{image\_00}, \texttt{image\_01} — 흑백 스테레오 (왼/오)
\item \texttt{oxts/data/*.txt} — 한 줄에 30개 값 (위경도/고도/자세/속도/가속도/각속도/\ldots)
  \begin{itemize}\small
  \item \texttt{values[14..16]} \(\to\) 가속도 (body frame, m/s\(^2\))
  \item \texttt{values[20..22]} \(\to\) 각속도 (body frame, rad/s)
  \item lat/lon/alt \(\to\) Mercator 투영 \(\to\) GT 위치
  \end{itemize}
\item LiDAR (\texttt{velodyne\_points/}) 와 tracklets 는 사용 안 함
\end{itemize}

\section{시스템 동작을 한 그림으로}

\begin{center}
\begin{tikzpicture}[
  >={Stealth[length=2.2mm]},
  every node/.style={font=\small},
  sensor/.style={draw, rectangle, rounded corners=2pt, minimum width=2.6cm, minimum height=0.95cm, align=center, fill=blue!6, inner sep=2pt},
  proc/.style={draw, rectangle, rounded corners=2pt, minimum width=3.4cm, minimum height=0.95cm, align=center, fill=green!8, inner sep=2pt},
  fuse/.style={draw, rectangle, rounded corners=2pt, minimum width=2.5cm, minimum height=1.4cm, align=center, fill=green!12, inner sep=2pt, very thick},
  output/.style={draw, rectangle, rounded corners=2pt, minimum width=2.8cm, minimum height=0.85cm, align=center, fill=orange!10, inner sep=2pt},
  gtbox/.style={draw, rectangle, rounded corners=2pt, minimum width=2.6cm, minimum height=0.95cm, align=center, fill=yellow!15, inner sep=2pt, dashed},
  arr/.style={->, thick},
  darr/.style={->, thick, dashed}
]
% ---- Layer 1: Sensors (column 1) ----
\node[sensor] (cam)  at (0, 2.4) {Stereo PNG \\[-0.1em] {\scriptsize image\_00 \& image\_01}};
\node[sensor] (imu)  at (0, 0.0) {OXTS IMU \\[-0.1em] {\scriptsize 10 Hz, sync}};
\node[gtbox]  (gt)   at (0,-2.4) {OXTS GPS (GT) \\[-0.1em] {\scriptsize 비교용, 추정에 미사용}};

% ---- Layer 2: Frontend (column 2) ----
\node[proc] (vo)   at (4.5, 2.4) {StereoTracker \\[-0.1em] {\scriptsize KLT $+$ ORB $+$ 삼각측량 $+$ PnP}};
\node[proc] (prop) at (4.5, 0.0) {ImuPropagator \\[-0.1em] {\scriptsize RK4 적분 $+$ 공분산 전파}};

% ---- Layer 3: EKF (column 3) ----
\node[fuse] (ekf) at (9.0, 1.2) {LcEkf \\[0.15em] {\scriptsize Kalman 보정} \\[-0.1em] {\scriptsize (15-state)}};

% ---- Layer 4: Outputs (column 4) ----
\node[output] (out1) at (12.8,  2.4) {trajectory\_tum.txt};
\node[output] (out2) at (12.8,  0.0) {ekf 위치 / 자세};
\node[output] (out3) at (12.8, -2.4) {metrics.txt \\[-0.1em] {\scriptsize ATE RMSE}};

% ---- Arrows: sensors -> processing ----
\draw[arr] (cam)  -- node[above, font=\scriptsize] {좌/우 이미지} (vo);
\draw[arr] (imu)  -- node[above, font=\scriptsize] {gyro, accel} (prop);

% ---- Arrows: processing -> EKF ----
\draw[arr] (vo.east)   -- node[above, sloped, font=\scriptsize] {p\_world\_cam0} (ekf.north west);
\draw[arr] (prop.east) -- node[below, sloped, font=\scriptsize] {p, v, R, P} (ekf.south west);

% ---- Arrows: EKF -> outputs ----
\draw[arr] (ekf.north east) -- (out1.west);
\draw[arr] (ekf.east)       -- (out2.west);

% ---- GT comparison path ----
\draw[darr] (gt.east) -- ++(7.5,0) |- (out3.west);
\node[font=\scriptsize] at (8.0,-2.7) {Umeyama 정합};
\draw[darr] (out2.south) -- (out3.north);

% ---- Group labels ----
\node[font=\footnotesize\bfseries, blue!50!black] at (0,3.5)  {입력 센서};
\node[font=\footnotesize\bfseries, green!50!black] at (4.5,3.5) {프론트엔드};
\node[font=\footnotesize\bfseries, green!30!black] at (9.0,3.5) {융합 (EKF)};
\node[font=\footnotesize\bfseries, orange!60!black] at (12.8,3.5) {출력};
\end{tikzpicture}
\end{center}

\begin{center}\footnotesize
\textit{실선 = 데이터 흐름, 점선 = GT 비교만 (추정 단계엔 사용 안 함)}
\end{center}

\section{처음부터 끝까지 무엇을 하나 — 5단계 1시간 30분}

\begin{center}\small
\begin{tabular}{rlll}
\toprule
단계 & 내용 & 시간 & 가이드 \\
\midrule
1 & WSL Ubuntu + apt 의존성 설치 & 30 분 & 운용 절차 \\
2 & KITTI 계정 → calib + 1 drive sync 다운 & 30 분 & 운용 절차 \\
3 & calib 3 종 → \(T_{\text{cam,imu}}\) 추출 → yaml & 15 분 & 캘리브레이션/YAML \\
4 & cmake + make → \texttt{build/run\_vio} & 5 분 & 빌드 \& 실행 \\
5 & 단일 시퀀스 실행 + plot 확인 & 1\textasciitilde{}5 분 & 실행/시각화 \\
\bottomrule
\end{tabular}
\end{center}

여기까지 \textbf{약 1시간 30분} 이면 첫 결과가 나온다. 그 후 확장: 더 긴 시퀀스, 파라미터 튜닝, 다중 시퀀스 비교.

\section{첫 결과 — ``정상''의 기준}

\texttt{drive\_0001} (11초, 108프레임) 으로 처음 돌리면 stdout 끝에:

\begin{lstlisting}[style=shell]
========================================
Frames processed : 108
ATE RMSE         : 156 cm
========================================
\end{lstlisting}

\textbf{ATE RMSE 1\textasciitilde{}3 m} 가 LC-EKF 정상 범위. 같은 시퀀스를 Tight MSCKF (OpenVINS 등) 로 돌리면 5\textasciitilde{}30 cm 까지 줄어든다 (5\textasciitilde{}30배 차이). 이게 LC-EKF 의 \textbf{본질적 천장}이다.

\texttt{trajectory\_plot.png} 의 왼쪽 패널 (XY top-down) 에서 추정(파란)과 GT(초록)이 \textbf{같은 모양}으로 겹쳐 보이면 정상. 약간의 평행이동/회전 오차는 LC-EKF 의 알려진 한계 — 글로벌 yaw 관측이 없어 절대 방향이 점점 어긋난다.

\section{용어집}

\begin{center}\small
\begin{longtable}{lp{0.7\textwidth}}
\toprule
용어 & 의미 \\
\midrule\endhead
VIO & Visual-Inertial Odometry. 카메라+IMU 위치 추적 \\
VO & Visual Odometry. 이미지만으로 카메라 위치 추적 \\
SLAM & Simultaneous Localization And Mapping. 위치 + 지도 동시 추정 \\
EKF & Extended Kalman Filter. 비선형 상태 추정의 표준 도구 \\
LC-EKF & Loosely-Coupled EKF. VO 의 출력(위치)을 측정값으로 받음 (이 프로젝트) \\
Tight MSCKF & 특징점 픽셀을 직접 측정값으로 받음 (OpenVINS, MSCKF) \\
Stereo & 좌/우 두 카메라로 깊이 추정 \\
Baseline & 두 카메라 사이 물리 거리 (KITTI: 53.7 cm) \\
Disparity & 같은 점이 좌\(\leftrightarrow\)우 이미지에서 픽셀 차이 \\
KLT & Kanade-Lucas-Tomasi 광류. 프레임간 특징점 추적 (작은 이동) \\
ORB & Oriented FAST + Rotated BRIEF. 회전·스케일 강건 디스크립터 (스테레오 매칭용) \\
FAST & Features from Accelerated Segment Test. 빠른 코너 검출기 \\
PnP & Perspective-n-Point. 2D-3D 대응에서 카메라 pose 계산 \\
RANSAC & RANdom SAmple Consensus. outlier 강건 추정 \\
IMU & Inertial Measurement Unit. 3축 가속도 + 3축 자이로 \\
OXTS & KITTI에 장착된 GPS+IMU 통합 유닛 (RT3003 모델) \\
GT & Ground Truth. 정답값 (KITTI 는 OXTS GPS+RTK) \\
ATE & Absolute Trajectory Error. 정합 후 잔차 RMSE \\
RMSE & Root Mean Square Error \\
TUM 포맷 & timestamp + xyz + 쿼터니언 한 줄씩의 trajectory 저장 형식 \\
Rectification & 스테레오 좌/우 이미지의 epipolar line 을 수평 스캔라인으로 만드는 변환 \\
Lever arm & IMU 와 카메라 위치 사이의 물리 벡터 (\(p_{IC}\)) \\
6-DoF pose & 3D 위치 + 3D 회전 = 6 자유도 자세 \\
RK4 & Runge-Kutta 4 차. 미분방정식 수치 적분기 \\
Umeyama & 두 점군을 강체 정합하는 SVD 기반 알고리즘 \\
\bottomrule
\end{longtable}
\end{center}


% =====================================================================
\chapter{수학적 배경 — 왜 그 수식인가}
% =====================================================================

이 장은 \textbf{각 수식이 왜 그렇게 쓰여 있는가} 를 설명한다. 도출 과정 자체보다는 \emph{설계 선택의 이유}에 초점이 있다.

% ---------------------------------------------------------------------
\section{회전 표현: SO(3) vs 오일러각 vs 쿼터니언}
% ---------------------------------------------------------------------

\subsection{왜 오일러각이 아닌가}

오일러각 \((\phi, \theta, \psi)\) 은 직관적이지만 두 가지 본질적 결함:
\begin{enumerate}
\item \textbf{Gimbal lock}: pitch \(\theta = \pm \pi/2\) 에서 yaw 와 roll 이 같은 자유도가 되어 한 자유도가 사라진다. 차량은 \(\pm \pi/2\) 까지 갈 일이 거의 없지만, 드론/로봇팔에선 빈번하다.
\item \textbf{회전 합성이 비선형}: \(R(\phi_1, \theta_1, \psi_1) \cdot R(\phi_2, \theta_2, \psi_2)\) 의 결과를 다시 오일러각으로 뽑으려면 비자명한 식이 필요하다. EKF 에서 \(R \leftarrow R \cdot R_\delta\) 같은 합성을 매 스텝 해야 하는데, 이 비선형성이 누적된다.
\end{enumerate}

\subsection{왜 쿼터니언만으로는 부족한가}

쿼터니언 \(q = (q_w, q_x, q_y, q_z)\) 은 4개 값으로 회전을 나타내며 gimbal lock 이 없다. 그러나:
\begin{itemize}
\item \textbf{유닛 제약 \(\|q\| = 1\)}: 4개 값이지만 자유도는 3. EKF 가 4D 상태로 쿼터니언을 추정하면 매 업데이트마다 정규화로 매니폴드 위로 끌어와야 하고, 공분산이 4D 인데 실제 자유도가 3D 라 \emph{수치적 비조건수}(numerical ill-conditioning) 가 생긴다.
\item \textbf{이중 표지(double cover)}: \(q\) 와 \(-q\) 는 같은 회전이라 부호 모호성이 있다.
\end{itemize}

\subsection{왜 SO(3) 와 error-state 가 정답인가}

\textbf{SO(3)} 는 \(3 \times 3\) 직교행렬 \(R\) 로서 \(R\T R = I, \det R = 1\) 인 3차원 \emph{매니폴드}이다. 그 \textbf{접공간(tangent space)} 은 3차원 Lie algebra \(\so(3)\) 으로, 작은 각도 벡터 \(\delta\theta \in \R^3\) 으로 표현된다.

이걸 활용한 \textbf{error-state EKF} 의 핵심 아이디어:
\begin{itemize}
\item \textbf{Nominal state}: 큰 회전을 \(R \in \SO(3)\) (\(3\times 3\) 행렬) 로 그대로 들고 다닌다.
\item \textbf{Error state}: 작은 회전 오차만 \(\delta\theta \in \R^3\) 로 추정한다 — 이 부분이 \emph{선형 EKF} 가 다룬다.
\item \textbf{합성}: 보정 시 \(R \leftarrow R \cdot \Exp(\delta\theta)\) 로 매니폴드 위에 정확히 머물게 한다.
\end{itemize}

이 분리로 얻는 것:
\begin{enumerate}
\item EKF 가 다루는 차원이 15 (4D 쿼터니언 시 16) 로 정확히 자유도와 일치 \(\to\) 공분산이 항상 well-conditioned.
\item \(R\) 자체는 매니폴드를 절대 벗어나지 않음 (정규화 후 cosmetics 만 필요).
\item Jacobian 이 단순해짐 (3D 작은 각도에 대한 미분).
\end{enumerate}

\subsection{핵심 SO(3) 연산}

\textbf{Hat 연산자} \([\cdot]_\times : \R^3 \to \so(3)\):
\begin{equation}
\skewop{\bm{v}} =
\begin{bmatrix}
0 & -v_3 & v_2 \\
v_3 & 0 & -v_1 \\
-v_2 & v_1 & 0
\end{bmatrix}
\end{equation}

성질: \(\skewop{\bm{a}} \bm{b} = \bm{a} \times \bm{b}\). 외적을 행렬 곱으로 바꾸는 도구.

\textbf{지수 사상 (Rodrigues 공식)} \(\Exp : \R^3 \to \SO(3)\):
\begin{equation}
\Exp(\bm{\omega}) = I + \frac{\sin \theta}{\theta} \skewop{\bm{\omega}} + \frac{1 - \cos \theta}{\theta^2} \skewop{\bm{\omega}}^2, \quad \theta = \|\bm{\omega}\|
\end{equation}

작은 \(\bm{\omega}\) 에서 \(\Exp(\bm{\omega}) \approx I + \skewop{\bm{\omega}}\). 즉 \emph{first-order Taylor 가 정확히 작은 회전} 이다.

\textbf{로그 사상} \(\Log : \SO(3) \to \R^3\): \(R\) 에서 회전축\(\times\)각도 벡터를 추출.

이들이 \texttt{include/core/so3.hpp} 에 구현돼 있고, 모든 회전 합성과 미분에서 일관되게 쓰인다.

% ---------------------------------------------------------------------
\section{Error-state EKF 의 15차원 상태}
% ---------------------------------------------------------------------

\subsection{Nominal state \(\to\) Error state 분해}

\textbf{Nominal state} (실제 추정치):
\begin{equation}
\bm{x} = (\bm{p}, \bm{v}, R, \bm{b}_g, \bm{b}_a) \in \R^3 \times \R^3 \times \SO(3) \times \R^3 \times \R^3
\end{equation}

\textbf{Error state} (작은 편차):
\begin{equation}
\delta\bm{x} = (\delta\bm{p}, \delta\bm{v}, \delta\bm{\theta}, \delta\bm{b}_g, \delta\bm{b}_a) \in \R^{15}
\end{equation}

합성 규칙:
\begin{align*}
\bm{p}_{\text{true}} &= \bm{p} + \delta\bm{p} \\
\bm{v}_{\text{true}} &= \bm{v} + \delta\bm{v} \\
R_{\text{true}} &= R \cdot \Exp(\delta\bm{\theta}) \\
\bm{b}_{g,\text{true}} &= \bm{b}_g + \delta\bm{b}_g \\
\bm{b}_{a,\text{true}} &= \bm{b}_a + \delta\bm{b}_a
\end{align*}

위치/속도/바이어스는 가법, 회전만 \(\Exp\) 합성.

\subsection{왜 IMU 측정 모델이 이렇게 생겼나}

IMU 가 출력하는 raw 신호:
\begin{align}
\tilde{\bm{\omega}} &= \bm{\omega} + \bm{b}_g + \bm{n}_g \\
\tilde{\bm{a}} &= R\T(\bm{a} - \bm{g}) + \bm{b}_a + \bm{n}_a
\end{align}

핵심 포인트:
\begin{itemize}
\item 가속도계는 \emph{비중력 가속도}만 잰다 (자유낙하 시 0 출력). 그래서 ``world frame 가속도 \(\bm{a}\)'' 에서 ``world frame 중력 \(\bm{g}\)'' 을 빼고 body frame 으로 회전 변환한 게 raw 출력.
\item 따라서 \emph{정지 상태} 의 IMU 가속도계는 \(R\T(-\bm{g})\) 를 출력 — 즉 ``중력의 반대 방향'' 을 가리킨다. 이게 초기 자세 추정의 기반이다.
\end{itemize}

본 코드의 초기화 (\texttt{run\_vio.cpp} 의 \texttt{initialize\_from\_imu}) 가 정확히 이 원리를 쓴다: 1초 동안 IMU 평균을 내고, 그 평균 가속도가 \(R\T \cdot (0,0,9.81)\T\) 이 되도록 \(R_0\) 를 계산한다.

\subsection{왜 15차원인가 (16 도 17 도 아니고)}

자유도 회계:
\begin{itemize}
\item 위치 3 + 속도 3 = 6
\item 회전 3 (SO(3) 의 차원)
\item 자이로 bias 3 + 가속도 bias 3 = 6
\item \textbf{합계 15}
\end{itemize}

쿼터니언으로 표현했다면 16D 가 되지만 자유도는 여전히 15. SO(3) 표현이 자유도와 정확히 일치한다.

% ---------------------------------------------------------------------
\section{IMU 운동학과 RK4 적분}
% ---------------------------------------------------------------------

\subsection{연속시간 운동학}

bias 보정된 IMU 측정값을 \(\bm{\omega}_c = \tilde{\bm{\omega}} - \bm{b}_g\), \(\bm{a}_c = \tilde{\bm{a}} - \bm{b}_a\) 로 두면:
\begin{align}
\dot{\bm{p}} &= \bm{v} \\
\dot{\bm{v}} &= R \bm{a}_c + \bm{g} \\
\dot{R} &= R \skewop{\bm{\omega}_c}
\end{align}

(bias 는 random walk 가정: \(\dot{\bm{b}}_g = \bm{n}_{bg}, \dot{\bm{b}}_a = \bm{n}_{ba}\))

\subsection{왜 RK4 인가}

가장 단순한 forward Euler:
\(\bm{p}_{k+1} = \bm{p}_k + dt \cdot \bm{v}_k\)
는 1차 정확도라 오차가 \(O(dt)\) 로 누적된다. IMU 가 10 Hz (\(dt = 0.1\) s) 만 되어도 1초 동안 가속이 변하면 무시 못 할 오차가 쌓인다.

RK4 (Runge-Kutta 4차) 는 한 스텝 안에서 4번 도함수를 평가해 \(O(dt^5)\) 정확도를 낸다 (지역 오차). 같은 \(dt = 0.1\) s 에서 forward Euler 대비 \(\sim 10^4\) 배 정확.

본 코드 (\texttt{imu\_propagator.cpp:40-66}) 의 \texttt{rk4\_step}:

\begin{lstlisting}[style=cpp,caption={RK4 한 스텝 — \texttt{src/ekf/imu\_propagator.cpp}}]
void ImuPropagator::rk4_step(double dt,
                              const Vector3d& gyro_c,
                              const Vector3d& accel_c) {
    State s0{p, v, R};
    auto k1 = state_derivative(s0, gyro_c, accel_c);

    State s1{p + 0.5*dt*k1.p,
             v + 0.5*dt*k1.v,
             R * so3::Exp(0.5*dt*gyro_c)};
    auto k2 = state_derivative(s1, gyro_c, accel_c);

    State s2{p + 0.5*dt*k2.p,
             v + 0.5*dt*k2.v,
             R * so3::Exp(0.5*dt*gyro_c)};
    auto k3 = state_derivative(s2, gyro_c, accel_c);

    State s3{p + dt*k3.p,
             v + dt*k3.v,
             R * so3::Exp(dt*gyro_c)};
    auto k4 = state_derivative(s3, gyro_c, accel_c);

    p = p + (dt/6.0) * (k1.p + 2*k2.p + 2*k3.p + k4.p);
    v = v + (dt/6.0) * (k1.v + 2*k2.v + 2*k3.v + k4.v);
    R = R * so3::Exp(dt * gyro_c);
    R = Eigen::Quaterniond(R).normalized().toRotationMatrix();
}
\end{lstlisting}

\textbf{한 가지 단순화}: 회전은 \(R \cdot \Exp(dt \, \bm{\omega}_c)\) 로 한 번에 적용한다 (RK4 수준의 회전은 Magnus expansion 이 필요한데, 차량 동역학에선 무시 가능).

% ---------------------------------------------------------------------
\section{공분산 전파: F, G, Q 행렬 유도}
% ---------------------------------------------------------------------

\subsection{Error-state ODE}

\(\delta\bm{x}\) 의 시간 도함수를 nominal state 주변에서 1차 선형화하면:
\begin{equation}
\dot{\delta\bm{x}} = F \delta\bm{x} + G \bm{n}
\end{equation}

여기서 \(\bm{n} = (\bm{n}_g, \bm{n}_a, \bm{n}_{bg}, \bm{n}_{ba}) \in \R^{12}\) 는 잡음 벡터.

\subsection{F 행렬 — 각 블록의 유도}

\textbf{상태 정의}: \(\delta\bm{x} = [\delta\bm{p},\; \delta\bm{v},\; \delta\bm{\theta},\; \delta\bm{b}_g,\; \delta\bm{b}_a] \in \R^{15}\).
Right-perturbation 규약 \(R_{\text{true}} = \hat{R}\,\Exp(\delta\bm{\theta})\) 를 채택한다.

\textbf{블록별 유도}

\medskip
\noindent\textbf{(1) \(\dot{\delta\bm{p}} = \delta\bm{v}\)} \hfill \(\Rightarrow F_{pp}=0,\; F_{pv}=I\)

위치의 시간 미분은 속도이므로 자명하다.

\medskip
\noindent\textbf{(2) \(\dot{\delta\bm{v}}\) 유도} \hfill \(\Rightarrow F_{v\theta}=-R\skewop{\bm{a}_c},\; F_{vb_a}=-R\)

진실 속도 미분과 nominal 속도 미분의 차:
\begin{align}
\dot{\bm{v}}_{\text{true}} &= R_{\text{true}}\,\bm{a}_{c,\text{true}} + \bm{g}
  = \hat{R}\,\Exp(\delta\bm{\theta})\,(\bm{a}_c - \delta\bm{b}_a + \bm{n}_a) + \bm{g} \notag\\
\dot{\hat{\bm{v}}} &= \hat{R}\,\bm{a}_c + \bm{g} \notag
\end{align}
\(\Exp(\delta\bm{\theta}) \approx I + \skewop{\delta\bm{\theta}}\) 로 1차 전개하면:
\begin{equation}
\delta\dot{\bm{v}} \approx \hat{R}\,\skewop{\delta\bm{\theta}}\,\bm{a}_c \;-\; \hat{R}\,\delta\bm{b}_a \;+\; \hat{R}\,\bm{n}_a
\end{equation}
\(\skewop{\delta\bm{\theta}}\,\bm{a}_c = \delta\bm{\theta} \times \bm{a}_c = -\skewop{\bm{a}_c}\,\delta\bm{\theta}\) 이므로:
\begin{equation}
\delta\dot{\bm{v}} = \underbrace{-R\skewop{\bm{a}_c}}_{F_{v\theta}}\,\delta\bm{\theta}
                  + \underbrace{(-R)}_{F_{vb_a}}\,\delta\bm{b}_a
                  + R\,\bm{n}_a
\end{equation}

\medskip
\noindent\textbf{(3) \(\dot{\delta\bm{\theta}}\) 유도} \hfill \(\Rightarrow F_{\theta\theta}=-\skewop{\bm{\omega}_c},\; F_{\theta b_g}=-R\)

회전 운동 방정식 \(\dot{R} = R\,\skewop{\bm{\omega}}\) 에서 right-perturbation으로 전개하면:
\begin{equation}
\frac{d}{dt}\bigl(\hat{R}\,\Exp(\delta\bm{\theta})\bigr)
  = \hat{R}\,\Exp(\delta\bm{\theta})\,\skewop{\bm{\omega}_c - \delta\bm{b}_g + \bm{n}_g}
\end{equation}
좌변의 product rule 과 \(\Exp(\delta\bm{\theta}) \approx I + \skewop{\delta\bm{\theta}}\) 를 적용하면:
\begin{equation}
\dot{\hat{R}}\,\Exp(\delta\bm{\theta}) + \hat{R}\,\Exp(\delta\bm{\theta})\,\skewop{\delta\dot{\bm{\theta}}}
\approx \hat{R}\,\bigl(I+\skewop{\delta\bm{\theta}}\bigr)\,\skewop{\bm{\omega}_c - \delta\bm{b}_g + \bm{n}_g}
\end{equation}
\(\dot{\hat{R}} = \hat{R}\,\skewop{\bm{\omega}_c}\) 를 대입하고 양변에 \(\hat{R}^{-1}\) 를 곱한 뒤 1차 항만 취하면:
\begin{equation}
\skewop{\delta\dot{\bm{\theta}}} \approx
  -\skewop{\bm{\omega}_c}\,\skewop{\delta\bm{\theta}}
  + \skewop{\delta\bm{\theta}}\,\skewop{\bm{\omega}_c}
  - \skewop{\delta\bm{b}_g} + \skewop{\bm{n}_g}
\end{equation}
교환자 \([\skewop{a},\skewop{b}] = \skewop{a \times b}\) 를 정리하면 벡터 형태로:
\begin{equation}
\delta\dot{\bm{\theta}} = \underbrace{-\skewop{\bm{\omega}_c}}_{F_{\theta\theta}}\,\delta\bm{\theta}
                        + \underbrace{(-R)}_{F_{\theta b_g}}\,\delta\bm{b}_g
                        + R\,\bm{n}_g
\end{equation}
\(F_{\theta b_g} = -R\) 인 이유: bias 는 body frame 잡음이고 \(\delta\bm{\theta}\) 는 world frame 표현이라 \(R\) 이 필요하다.

\medskip
\noindent\textbf{(4) bias 행}: random walk 모델 \(\dot{\bm{b}}_{g/a} = \bm{n}_{bg/ba}\) 에서 F 의 bias 행은 모두 0.

\medskip
이를 합치면 \(15{\times}15\) 연속 전이 행렬:

\begin{equation}
\boxed{
F = \begin{bmatrix}
\bm{0} & I & \bm{0} & \bm{0} & \bm{0} \\[3pt]
\bm{0} & \bm{0} & -R\skewop{\bm{a}_c} & \bm{0} & -R \\[3pt]
\bm{0} & \bm{0} & -\skewop{\bm{\omega}_c} & -R & \bm{0} \\[3pt]
\bm{0} & \bm{0} & \bm{0} & \bm{0} & \bm{0} \\[3pt]
\bm{0} & \bm{0} & \bm{0} & \bm{0} & \bm{0}
\end{bmatrix}
}
\quad
\begin{array}{l}
\leftarrow \delta\bm{p} \\
\leftarrow \delta\bm{v} \\
\leftarrow \delta\bm{\theta} \\
\leftarrow \delta\bm{b}_g \\
\leftarrow \delta\bm{b}_a
\end{array}
\end{equation}

코드 매핑 (\texttt{imu\_propagator.cpp:93-99}):
\begin{center}\small
\begin{tabular}{lll}
\toprule
블록 & 코드 & 의미 \\
\midrule
\(F[0,3]=I\)             & \texttt{F.block<3,3>(0,3)  = I}                   & \(\dot{\delta\bm{p}}=\delta\bm{v}\) \\
\(F[3,6]=-R[\bm{a}_c]_\times\) & \texttt{F.block<3,3>(3,6)  = -R*hat(accel\_c)}    & 자세 오차 \(\to\) 속도 미분 \\
\(F[3,12]=-R\)           & \texttt{F.block<3,3>(3,12) = -R}                  & accel bias \(\to\) 속도 미분 \\
\(F[6,6]=-[\bm{\omega}_c]_\times\) & \texttt{F.block<3,3>(6,6)  = -hat(gyro\_c)}  & 자세 error 자체 dynamics \\
\(F[6,9]=-R\)            & \texttt{F.block<3,3>(6,9)  = -R}                  & gyro bias \(\to\) 자세 미분 \\
\bottomrule
\end{tabular}
\end{center}

\subsection{G 행렬 — 잡음 입력 매핑}

잡음 벡터 \(\bm{n} = (\bm{n}_g,\;\bm{n}_a,\;\bm{n}_{bg},\;\bm{n}_{ba}) \in \R^{12}\) 가 error state 에 어떻게 들어오는지 위 유도에서 직접 읽을 수 있다:

\begin{itemize}
\item \(\delta\dot{\bm{v}}\) 항의 \(R\,\bm{n}_a\): accel 잡음은 body frame (\(\bm{n}_a\) 단위: \(\text{m/s}^2/\sqrt{\text{Hz}}\)), 하지만 \(\delta\bm{v}\) 는 world frame → \(R\) 로 회전 필요.
\item \(\delta\dot{\bm{\theta}}\) 항의 \(R\,\bm{n}_g\): gyro 잡음은 body frame (\(\bm{n}_g\) 단위: \(\text{rad/s}/\sqrt{\text{Hz}}\)), \(\delta\bm{\theta}\) 는 world frame → 마찬가지로 \(R\) 필요.
\item bias random walk \(\bm{n}_{bg},\bm{n}_{ba}\) 는 각 bias 에 직접 더해지므로 \(I\).
\end{itemize}

따라서 \(G \in \R^{15 \times 12}\):

\begin{equation}
\boxed{
G = \begin{bmatrix}
\bm{0} & \bm{0} & \bm{0} & \bm{0} \\[3pt]
\bm{0} & R      & \bm{0} & \bm{0} \\[3pt]
R      & \bm{0} & \bm{0} & \bm{0} \\[3pt]
\bm{0} & \bm{0} & I      & \bm{0} \\[3pt]
\bm{0} & \bm{0} & \bm{0} & I
\end{bmatrix}
}
\quad
\begin{array}{l}
\leftarrow \delta\bm{p} \\
\leftarrow \delta\bm{v} \\
\leftarrow \delta\bm{\theta} \\
\leftarrow \delta\bm{b}_g \\
\leftarrow \delta\bm{b}_a
\end{array}
\qquad
\begin{array}{l}
\uparrow\,\bm{n}_g \\
\uparrow\,\bm{n}_a \\
\uparrow\,\bm{n}_{bg} \\
\uparrow\,\bm{n}_{ba}
\end{array}
\end{equation}

코드 매핑 (\texttt{imu\_propagator.cpp:106-110}):
\begin{center}\small
\begin{tabular}{lll}
\toprule
블록 (행,열 인덱스) & 코드 & 의미 \\
\midrule
\(G[3:6,\; 3:6]=R\)   & \texttt{G.block<3,3>(3,3)  = R} & accel 잡음 \(\to\) \(\delta\bm{v}\) \\
\(G[6:9,\; 0:3]=R\)   & \texttt{G.block<3,3>(6,0)  = R} & gyro 잡음 \(\to\) \(\delta\bm{\theta}\) \\
\(G[9:12,\; 6:9]=I\)  & \texttt{G.block<3,3>(9,6)  = I} & gyro walk \(\to\) \(\delta\bm{b}_g\) \\
\(G[12:15,\;9:12]=I\) & \texttt{G.block<3,3>(12,9) = I} & accel walk \(\to\) \(\delta\bm{b}_a\) \\
\bottomrule
\end{tabular}
\end{center}

\subsection{\(Q_c\) 행렬 — 연속시간 잡음 공분산}

잡음 모델은 white noise 밀도: \(\bm{n}_g \sim \mathcal{N}(\bm{0},\,\sigma_g^2 I)\) 등 (단위 \(\sigma/\sqrt{\text{Hz}}\), 즉 PSD 형태). 12$\times$12 대각 행렬:

\begin{equation}
Q_c = \diag\!\bigl(\underbrace{\sigma_g^2 I_3}_{\bm{n}_g},\;
                   \underbrace{\sigma_a^2 I_3}_{\bm{n}_a},\;
                   \underbrace{\sigma_{bg}^2 I_3}_{\bm{n}_{bg}},\;
                   \underbrace{\sigma_{ba}^2 I_3}_{\bm{n}_{ba}}\bigr)
     \in \R^{12\times 12}
\end{equation}

각 \(\sigma\) 는 YAML 의 \texttt{imu\_noise} 섹션에서 읽히며 (\texttt{imu\_propagator.cpp:112-119}):

\begin{center}\small
\begin{tabular}{llll}
\toprule
파라미터 & YAML 키 & 단위 & \(Q_c\) 대각 블록 \\
\midrule
\(\sigma_g\)  & \texttt{gyro\_noise}  & \(\text{rad/s}/\sqrt{\text{Hz}}\) & \(\sigma_g^2 I_3\)  \\
\(\sigma_a\)  & \texttt{accel\_noise} & \(\text{m/s}^2/\sqrt{\text{Hz}}\) & \(\sigma_a^2 I_3\)  \\
\(\sigma_{bg}\) & \texttt{gyro\_walk}  & \(\text{rad/s}^2/\sqrt{\text{Hz}}\) & \(\sigma_{bg}^2 I_3\) \\
\(\sigma_{ba}\) & \texttt{accel\_walk} & \(\text{m/s}^3/\sqrt{\text{Hz}}\) & \(\sigma_{ba}^2 I_3\) \\
\bottomrule
\end{tabular}
\end{center}

\subsection{이산화: \(\Phi\) 와 \(Q_d\)}

연속 ODE 를 dt 동안 적분 (\texttt{imu\_propagator.cpp:101-124}):

\begin{align}
\Phi &\approx I + F \, dt \quad \text{(1차 Euler 근사)} \\
Q_d &= G \, Q_c \, G\T \cdot dt \quad \text{(잡음 공분산 이산화)}
\end{align}

\textbf{왜 \(dt\) 를 곱하나}: \(Q_c\) 는 스펙트럼 밀도 (단위 \(\sigma^2/\text{Hz}\)) 이므로, 이를 \(dt\) 구간에서 적분하면 \(\sigma^2 \cdot dt\) 가 분산이 된다. 구체적으로 \(Q_d\) 의 \((i,j)\) 블록은:
\begin{equation}
[Q_d]_{ij} = \sigma_i^2 \cdot [G G\T]_{ij} \cdot dt
\end{equation}

\textbf{전체 공분산 전파} (\texttt{imu\_propagator.cpp:124}):
\begin{equation}
\boxed{P_{k+1} = \Phi \, P_k \, \Phi\T + Q_d}
\end{equation}

코드에서는 \texttt{Phi = I + F*dt}, \texttt{Qd = G*Qc*G.transpose()*dt} 로 계산한 뒤 한 줄로 적용:
\begin{lstlisting}[style=cpp, caption={imu\_propagator.cpp:124}]
P = Phi * P * Phi.transpose() + Qd;
\end{lstlisting}

% ---------------------------------------------------------------------
\section{측정 모델과 Kalman 보정}
% ---------------------------------------------------------------------

\subsection{LC-EKF 측정 모델}

VO 가 출력한 cam0 의 world frame 위치 \(\bm{p}_{\text{cam0}}^W\) 를 측정값으로:
\begin{equation}
\bm{z} = \bm{p}_{\text{cam0}}^W = R^W_I \, \bm{p}_{IC} + \bm{p}^W_I + \bm{n}_v
\end{equation}

여기서:
\begin{itemize}
\item \(\bm{p}^W_I\) = IMU 의 world frame 위치 (EKF state).
\item \(R^W_I\) = world \(\leftarrow\) body 회전 (EKF state).
\item \(\bm{p}_{IC}\) = body frame 에서 본 cam0 원점 (\textbf{lever arm}, calibration 으로 고정).
\item \(\bm{n}_v\) = VO 측정 잡음, \(\sim \mathcal{N}(\bm{0}, \sigma_{\text{vo}}^2 I_3)\).
\end{itemize}

\subsection{H 행렬 유도 — \texttt{lc\_ekf.cpp:38-43}}

예측 측정값 \(h(\bm{x}) = R^W_I\,\bm{p}_{IC} + \bm{p}^W_I\) 를 error state 로 선형화한다. Right-perturbation:
\begin{equation}
R_{\text{true}} = \hat{R}^W_I\,\Exp(\delta\bm{\theta}), \qquad
\bm{p}_{\text{true}} = \hat{\bm{p}}^W_I + \delta\bm{p}
\end{equation}

예측값에서의 편차:
\begin{align}
\delta h &= h\bigl(\hat{\bm{x}} \oplus \delta\bm{x}\bigr) - h(\hat{\bm{x}}) \notag\\
&= \hat{R}\,\Exp(\delta\bm{\theta})\,\bm{p}_{IC} + \hat{\bm{p}} + \delta\bm{p}
   \;-\; \bigl(\hat{R}\,\bm{p}_{IC} + \hat{\bm{p}}\bigr) \notag\\
&\approx \hat{R}\,(I + \skewop{\delta\bm{\theta}})\,\bm{p}_{IC} - \hat{R}\,\bm{p}_{IC} + \delta\bm{p} \notag\\
&= \hat{R}\,\skewop{\delta\bm{\theta}}\,\bm{p}_{IC} + \delta\bm{p}
\end{align}

핵심 단계: \(\skewop{\delta\bm{\theta}}\,\bm{p}_{IC} = \delta\bm{\theta} \times \bm{p}_{IC} = -\bm{p}_{IC}\times\delta\bm{\theta} = -\skewop{\bm{p}_{IC}}\,\delta\bm{\theta}\)

따라서:
\begin{equation}
\delta h = \delta\bm{p} \;-\; \hat{R}\,\skewop{\bm{p}_{IC}}\,\delta\bm{\theta}
\end{equation}

각 편미분:
\begin{align}
\frac{\partial h}{\partial \delta\bm{p}} &= I_3 \\
\frac{\partial h}{\partial \delta\bm{\theta}} &= -R^W_I\,[\bm{p}_{IC}]_\times \\
\frac{\partial h}{\partial \delta\bm{v}} &= 0, \quad
\frac{\partial h}{\partial \delta\bm{b}_g} = 0, \quad
\frac{\partial h}{\partial \delta\bm{b}_a} = 0
\end{align}

\begin{kbox}
\textbf{주의}: \(-R[\bm{p}_{IC}]_\times \neq [R\bm{p}_{IC}]_\times\) (일반적으로 같지 않다).
표준 항등식 \([R\bm{v}]_\times = R\,[\bm{v}]_\times\,R^\top\) 에 의해
\(-R\,[\bm{p}_{IC}]_\times = -[R\bm{p}_{IC}]_\times\,R \neq [R\bm{p}_{IC}]_\times\).
올바른 야코비안은 \(-R\,[\bm{p}_{IC}]_\times\) 이다.
\end{kbox}

따라서:
\begin{equation}
\boxed{H = \begin{bmatrix} I_3 & \bm{0} & -R^W_I\,[\bm{p}_{IC}]_\times & \bm{0} & \bm{0} \end{bmatrix} \in \R^{3 \times 15}}
\end{equation}

코드 (\texttt{lc\_ekf.cpp:41-43}):
\begin{lstlisting}[style=cpp]
H.block<3,3>(0,0) = Matrix3d::Identity();   // dh/d(delta_p) = I
H.block<3,3>(0,6) = -R_wI * so3::hat(p_IC_); // dh/d(delta_theta) = -R*[p_IC]x
\end{lstlisting}

\textbf{lever arm 영향}: KITTI 에서 \(\bm{p}_{IC} \approx (-0.3,\; 0.7,\; -1.1)^{\!\top}\) m (약 1.3 m). 이 값이 0 이면 \(H_\theta = 0\) 이 되고, 클수록 자세 오차가 위치 측정에 더 크게 투영된다. 야코비안을 정확히 계산하는 것이 EKF 수렴에 중요하다.

\subsection{Kalman gain 과 상태 보정}

표준 EKF 업데이트 (\texttt{lc\_ekf.cpp:46-56}):
\begin{align}
S &= H P H\T + R_{\text{noise}} & \text{(innovation covariance)} \\
K &= P H\T S^{-1} & \text{(Kalman gain)} \\
\delta\bm{x} &= K (\bm{z} - h(\bm{x})) & \text{(error correction)}
\end{align}

해석:
\begin{itemize}
\item \(R_{\text{noise}} = \sigma_{\text{vo}}^2 I_3\) 가 작을수록 (VO 신뢰) \(K\) 가 커진다 \(\to\) 보정량 큼.
\item \(P\) 가 클수록 (IMU 의 추정 불확실) \(K\) 가 커진다 \(\to\) 측정값에 더 의존.
\item 둘의 비율이 \(K\) 의 핵심.
\end{itemize}

\subsection{왜 Joseph form 인가}

전통적 공분산 업데이트 \(P^+ = (I - KH) P\) 는 수학적으로 맞지만 \textbf{수치적으로 위험}:
\begin{itemize}
\item 부동소수 오차로 \(P^+\) 가 비대칭이 되거나 음의 eigenvalue 가 생길 수 있음.
\item 특히 측정값이 매우 신뢰될 때 (\(R_{\text{noise}} \to 0\)) 더 심해짐.
\end{itemize}

Joseph form (\texttt{lc\_ekf.cpp:60-61}):
\begin{equation}
\boxed{P^+ = (I - KH) P (I - KH)\T + K R_{\text{noise}} K\T}
\end{equation}

이 형태는:
\begin{itemize}
\item \textbf{대칭성 보장}: 두 항 모두 \(A B A\T\) 꼴이라 자동으로 대칭.
\item \textbf{양의 정부호 보장}: \(P, R_{\text{noise}}\) 가 PSD 면 결과도 PSD.
\item 추가 비용은 \(O(n^3)\) 한 번 더 (15D 에선 무시 가능).
\end{itemize}

% ---------------------------------------------------------------------
\section{Stereo Triangulation 과 PnP}
% ---------------------------------------------------------------------

\subsection{왜 베이스라인이 미터 스케일을 결정하나}

Pinhole + rectified stereo 모델에서 한 점이 좌측 픽셀 \((u_L, v)\), 우측 \((u_R, v)\) 에 보일 때:
\begin{equation}
Z = \frac{f_x \cdot b}{u_L - u_R}, \quad X = \frac{(u_L - c_x) \, Z}{f_x}, \quad Y = \frac{(v - c_y) \, Z}{f_y}
\end{equation}

\(b\) 는 미터 단위 베이스라인. 모든 깊이 \(Z\) 가 미터로 자동 결정된다. KITTI 의 \(b = 0.537\) m 가 calibration 시 레이저로 측정되어 yaml 의 \(T_{\text{cam,imu}}\) 에 박혀 있다.

코드에서는 OpenCV \texttt{cv::triangulatePoints} 가 \(P_0, P_1\) projection matrix 를 받아 4D 동차좌표를 출력한다 (\texttt{stereo\_tracker.cpp}).

\subsection{PnP 가 풀어내는 것}

\(N\) 개의 2D-3D 대응 \((u_i, v_i, X_i, Y_i, Z_i)\) 에서 카메라 pose \((R, \bm{t})\) 를 풀어내는 문제. 최소 \(N = 4\) (P3P 는 3 + ambiguity, EPnP/UPnP 는 4 이상 권장).

본 코드는 \texttt{cv::solvePnPRansac} 사용 — RANSAC outlier 제거 후 EPnP + reprojection 비선형 최적화.

\textbf{왜 단순 LSQ 가 아닌 RANSAC?} KLT/ORB 매칭에서 outlier (잘못된 대응) 가 5\textasciitilde{}30\% 발생. 단순 최소제곱은 outlier 한 개에 휘둘리지만 RANSAC 은 ``inlier 다수결'' 로 강건하다.

% ---------------------------------------------------------------------
\section{ATE 계산: Umeyama 정합}
% ---------------------------------------------------------------------

추정 trajectory \(\{\bm{e}_i\}\) 와 GT \(\{\bm{g}_i\}\) 간 \textbf{절대 trajectory 오차}:

\begin{kbox}
\textbf{왜 단순 \(\|\bm{e}_i - \bm{g}_i\|\) 가 아닌가?}\\[0.3em]
EKF 가 임의 시점을 원점으로 잡아 (\texttt{first IMU sample = (0,0,0)}) 시작하면 GT 와 \emph{전역 SE(3) 변환} 만큼 어긋난다. 그 변환을 ``드리프트'' 로 카운트하면 안 된다.
\end{kbox}

\textbf{Umeyama (1991)} 알고리즘 — SVD 로 두 점군의 최적 강체정합을 구함:
\begin{equation}
\min_{R, \bm{t}} \sum_i \| (R \bm{e}_i + \bm{t}) - \bm{g}_i \|^2
\end{equation}

\(\to\) \(\Sigma = \frac{1}{N} \sum_i (\bm{e}_i - \bar{\bm{e}})(\bm{g}_i - \bar{\bm{g}})\T\) 의 SVD \(U \Sigma V\T\) 로:
\begin{align}
R^* &= V \diag(1, 1, \det(VU\T)) U\T \\
\bm{t}^* &= \bar{\bm{g}} - R^* \bar{\bm{e}}
\end{align}

본 코드 (\texttt{run\_vio.cpp:193}) 는 Eigen 의 \texttt{Eigen::umeyama(\ldots, false)} 호출 — 두 번째 인자 \texttt{false} 는 ``스케일 고정'' 의미. 우리는 스테레오로 미터 스케일이 이미 맞으니 회전+평행이동만 정합.

\textbf{ATE RMSE}:
\begin{equation}
\text{ATE}_{\text{RMSE}} = \sqrt{\frac{1}{N} \sum_i \| (R^* \bm{e}_i + \bm{t}^*) - \bm{g}_i \|^2}
\end{equation}


% =====================================================================
\chapter{시스템 아키텍처}
% =====================================================================

\section{전체 흐름}

본 프로젝트는 ROS 없이 단일 C++ 실행파일로 구성된다. 입력 \(\to\) 처리 \(\to\) 출력의 큰 그림:

\begin{center}
\begin{tikzpicture}[
  >={Stealth[length=2mm]},
  every node/.style={font=\footnotesize},
  io/.style={draw, rectangle, rounded corners=2pt, fill=blue!8, minimum width=1.9cm, minimum height=0.75cm, align=center, inner sep=2pt},
  proc/.style={draw, rectangle, rounded corners=2pt, fill=green!8, minimum width=1.9cm, minimum height=0.75cm, align=center, inner sep=2pt},
  fuse/.style={draw, rectangle, rounded corners=2pt, fill=green!12, minimum width=1.9cm, minimum height=0.95cm, align=center, inner sep=2pt, very thick},
  outbox/.style={draw, rectangle, rounded corners=2pt, fill=orange!10, minimum width=2.1cm, minimum height=0.75cm, align=center, inner sep=2pt},
  arr/.style={->, thick}
]
% ---- Column x=0: inputs ----
\node[io]  (yaml) at (0, 3.0) {YAML config};
\node[io]  (data) at (0, 1.5) {KITTI data \\[-0.1em] {\tiny stereo PNG + OXTS}};
\node[io]  (gt)   at (0,-1.2) {GT \\[-0.1em] {\tiny OXTS GPS}};

% ---- Column x=3: reader ----
\node[proc] (rd)   at (3.2, 2.25) {KittiRawReader};

% ---- Column x=6.4: frontend ----
\node[proc] (track) at (6.4, 3.0) {StereoTracker};
\node[proc] (prop)  at (6.4, 1.5) {ImuPropagator};

% ---- Column x=9.6: EKF fusion ----
\node[fuse] (ekf)   at (9.6, 2.25) {LcEkf};

% ---- Column x=13: outputs and ATE ----
\node[outbox] (out1) at (13.2, 3.0) {trajectory\_tum.txt};
\node[proc]   (eval) at (13.2, 0.6) {compute\_ate};
\node[outbox] (out2) at (13.2,-1.2) {metrics.txt};

% ---- Arrows: inputs -> reader ----
\draw[arr] (yaml.east) -- ($(rd.west)+(0,0.25)$);
\draw[arr] (data.east) -- ($(rd.west)+(0,-0.15)$);

% ---- Reader -> frontend ----
\draw[arr] (rd.east) -- ($(track.west)$);
\draw[arr] (rd.east) -- ($(prop.west)$);

% ---- Frontend -> EKF ----
\draw[arr] (track.east) -- ($(ekf.west)+(0,0.25)$);
\draw[arr] (prop.east)  -- ($(ekf.west)+(0,-0.25)$);

% ---- EKF -> outputs (no crossing) ----
\draw[arr] (ekf.east) -- ($(out1.west)$);
\draw[arr] (ekf.south) |- ($(eval.west)+(0,0.15)$);

% ---- GT -> compute_ate (cleanly along bottom) ----
\draw[arr] (gt.east) -| ($(eval.west)+(0,-0.15)$);

% ---- compute_ate -> metrics ----
\draw[arr] (eval.south) -- (out2.north);

% ---- Group labels ----
\node[font=\scriptsize\bfseries, blue!50!black]   at (0,    3.9) {입력};
\node[font=\scriptsize\bfseries, green!50!black]  at (3.2,  3.9) {로더};
\node[font=\scriptsize\bfseries, green!50!black]  at (6.4,  3.9) {프론트엔드};
\node[font=\scriptsize\bfseries, green!30!black]  at (9.6,  3.9) {융합};
\node[font=\scriptsize\bfseries, orange!60!black] at (13.2, 3.9) {출력 / 평가};
\end{tikzpicture}
\end{center}

\section{발표용 코어 모듈 요약}

PPT 에서는 전체 코드를 모두 설명하기보다, 아래 3개 코어 모듈이 어떤 정보를 주고받는지 먼저 보여주는 편이 가장 명확하다.

\begin{kbox}
\[
\begin{aligned}
\texttt{StereoTracker} &\longrightarrow z_k = \bm{p}_{C_k}^{W} \\
\texttt{ImuPropagator} &\longrightarrow \bm{x}_{k+1}^{-},\ \bm{P}_{k+1}^{-}
\end{aligned}
\qquad \Longrightarrow \qquad
\texttt{LcEkf}
\]
\[
\bm{x} = \{\bm{p}, \bm{v}, R, \bm{b}_g, \bm{b}_a\}, \qquad
\delta\bm{x} =
\begin{bmatrix}
\delta\bm{p} & \delta\bm{v} & \delta\bm{\theta} & \delta\bm{b}_g & \delta\bm{b}_a
\end{bmatrix}^{\!\top}
\in \R^{15}
\]
\end{kbox}

\begin{center}\small
\begin{tabularx}{\textwidth}{p{0.20\textwidth}p{0.28\textwidth}X}
\toprule
모듈 & 발표용 한 줄 & PPT 에 넣을 핵심 수식 \\
\midrule
\texttt{StereoTracker} &
스테레오 이미지에서 sparse VO 위치 측정 생성 &
\[
\begin{aligned}
\min_{R_{CW},\,\bm{t}_{CW}}
\sum_i
\left\|
\bm{u}_i -
\pi\!\left(K(R_{CW}\bm{P}_i^W+\bm{t}_{CW})\right)
\right\|^2, \\
z_k &= \bm{p}_{C_k}^{W}
\end{aligned}
\]
\\
\addlinespace
\texttt{ImuPropagator} &
IMU 로 nominal state 와 covariance 예측 &
\[
\begin{aligned}
\dot{\bm{p}} &= \bm{v}, \\
\dot{\bm{v}} &= R(\bm{a}_m-\bm{b}_a)+\bm{g}, \\
\dot{R} &= R\skewop{\bm{\omega}_m-\bm{b}_g}
\end{aligned}
\]
\[
\begin{aligned}
\bm{P}_{k+1}^{-} &=
\Phi\bm{P}_{k}^{+}\Phi\T+\bm{G}\bm{Q}\bm{G}\T, \\
\Phi \approx I + F\Delta t
\end{aligned}
\]
\\
\addlinespace
\texttt{LcEkf} &
VO 위치 측정으로 IMU drift 보정 &
\[
\begin{aligned}
h(\bm{x}) &= \bm{p}_{I}^{W}+R_{I}^{W}\bm{p}_{IC}, \\
\bm{r} &= z-h(\bm{x})
\end{aligned}
\]
\[
H = \bigl[\;I_3 \;\;\bm{0}\;\; {-R^W_I[\bm{p}_{IC}]_\times} \;\;\bm{0}\;\;\bm{0}\;\bigr]
\]
\[
\begin{aligned}
K &= PH\T(HPH\T+R_n)^{-1}, \\
\delta\bm{x} &= K\bm{r}
\end{aligned}
\]
\\
\bottomrule
\end{tabularx}
\end{center}

\begin{notebox}
발표의 핵심 메시지: 이 시스템은 \textbf{Tight VIO} 가 아니라 \textbf{Loosely-Coupled VIO} 다.
즉 feature residual 을 EKF 에 직접 넣지 않고, \texttt{StereoTracker} 가 만든 camera position 을 EKF measurement 로 사용한다.
\end{notebox}

\section{모듈 책임 분리}

\begin{center}\small
\begin{tabular}{p{0.22\textwidth}p{0.22\textwidth}p{0.50\textwidth}}
\toprule
모듈 & 위치 & 책임 \\
\midrule
\texttt{KittiRawReader} & \texttt{src/io/kitti\_raw\_reader.cpp} &
디스크에서 KITTI raw 디렉터리 파싱. IMU/CAM/GT 를 \texttt{std::vector} 로 적재. \\
\addlinespace
\texttt{StereoTracker} & \texttt{src/frontend/stereo\_tracker.cpp} &
KLT + ORB + 삼각측량 + PnP 로 cam0 의 6-DoF pose 출력. 매니폴드 ``world = first cam0 frame''. \\
\addlinespace
\texttt{ImuPropagator} & \texttt{src/ekf/imu\_propagator.cpp} &
IMU 측정값으로 nominal state 와 공분산 \(P\) 전파 (RK4 + 1차 Euler). \\
\addlinespace
\texttt{LcEkf} & \texttt{src/ekf/lc\_ekf.cpp} &
\texttt{ImuPropagator} 를 wrap 하고 VO 측정으로 보정. Joseph form 공분산 업데이트. \\
\addlinespace
\texttt{run\_vio.cpp} & \texttt{apps/run\_vio.cpp} &
메인 루프: yaml 로드 → reader 생성 → 초기화 → IMU/Cam interleave → 결과 저장. \\
\bottomrule
\end{tabular}
\end{center}

\section{데이터 구조}

\texttt{include/core/types.hpp:1-37}:

\begin{lstlisting}[style=cpp,caption={핵심 struct}]
struct ImuData {
    double timestamp;       // seconds
    Eigen::Vector3d gyro;   // rad/s, body (IMU) frame
    Eigen::Vector3d accel;  // m/s^2, body (IMU) frame
};

struct CamData {
    double timestamp;
    std::string img_l;      // 좌측 png 절대경로
    std::string img_r;      // 우측 png 절대경로
};

struct GtData {
    double timestamp;
    Eigen::Vector3d p;      // world frame 위치
    Eigen::Quaterniond q;   // world ← body 쿼터니언
    Eigen::Vector3d v;      // world frame 속도
};

struct CameraParams {
    double fx, fy, cx, cy;          // intrinsics
    double k1, k2, p1, p2;          // radtan distortion
    Eigen::Matrix4d T_cam_imu;      // IMU → cam (4x4)
    int width, height;
};

struct ImuNoiseParams {
    double gyro_noise   = 1.6968e-4;  // sigma_g  [rad/s/sqrt(Hz)]
    double accel_noise  = 2.0000e-3;  // sigma_a
    double gyro_walk    = 1.9393e-5;  // sigma_bg
    double accel_walk   = 3.0000e-3;  // sigma_ba
};
\end{lstlisting}

\section{좌표계 5가지}

\begin{enumerate}
\item \textbf{World (W)}: 첫 IMU 샘플 시점에서의 IMU body frame. 원점 \((0,0,0)\), 중력은 \(-Z\).
\item \textbf{Body / IMU (I)}: 차량 IMU 좌표계. forward-X / left-Y / up-Z.
\item \textbf{Cam0 / Cam1 (C)}: 카메라 좌표계. right-X / down-Y / forward-Z (OpenCV 표준).
\item \textbf{Velodyne (V)}: 본 프로젝트 안 씀. \(T_{\text{cam,velo}}\) 만 calibration 추출 시 잠시 사용.
\item \textbf{VO local}: \texttt{StereoTracker} 가 첫 cam0 프레임을 원점 \((0,0,0)\) 으로 내부적으로 들고 다님. World 와 정렬하려면 \(R_{WC_0}, \bm{t}_{WC_0}\) 변환 필요.
\end{enumerate}

\textbf{Lever arm}: \(\bm{p}_{IC} = T_{IC} [\bm{0}; 1]\) 의 처음 3개 성분. IMU 좌표계에서 본 cam0 원점 위치. EKF 측정 모델의 핵심.


% =====================================================================
\chapter{모듈별 코드 분석}
% =====================================================================

이 장은 3개 코어 모듈과 2개 운용 모듈을 \emph{줄 단위로} 분석하고 앞 장의 수식과 1:1 매핑한다.

% ---------------------------------------------------------------------
\section{ImuPropagator — IMU 전파}
% ---------------------------------------------------------------------

\subsection{생성자: 초기 분산 설정}

\begin{lstlisting}[style=cpp,caption={\texttt{src/ekf/imu\_propagator.cpp:9-25}}]
ImuPropagator::ImuPropagator(const Vector3d& p0, const Vector3d& v0,
                              const Matrix3d& R0, const Vector3d& bg0,
                              const Vector3d& ba0,
                              const ImuNoiseParams& noise,
                              const Vector3d& g_world)
    : p(p0), v(v0), R(R0), bg(bg0), ba(ba0), noise_(noise), g_(g_world)
{
    P.setZero();
    P.block<3,3>(0,0)   = 1e-4 * Matrix3d::Identity();  // position
    P.block<3,3>(3,3)   = 1e-2 * Matrix3d::Identity();  // velocity
    P.block<3,3>(6,6)   = 1e-4 * Matrix3d::Identity();  // attitude
    P.block<3,3>(9,9)   = 1e-6 * Matrix3d::Identity();  // gyro bias
    P.block<3,3>(12,12) = 1e-4 * Matrix3d::Identity();  // accel bias
}
\end{lstlisting}

\textbf{왜 이 값들?}
\begin{itemize}
\item \(1e^{-4}\) m\(^2\) (위치): \(\sigma = 1\) cm. 첫 샘플은 ``\(0\) 에 가깝다''고 강하게 믿음.
\item \(1e^{-2}\) m\(^2\)/s\(^2\) (속도): \(\sigma = 0.1\) m/s. 차량이 정지일 가능성 ``상당히'' 믿음.
\item \(1e^{-4}\) (자세, rad\(^2\)): \(\sigma \approx 0.01\) rad \(\approx 0.6°\). 정지 시 가속도계로 자세를 잡았으니 신뢰 높음.
\item \(1e^{-6}\) (gyro bias): \(\sigma = 0.001\) rad/s. 1초 평균 후 잔차 추정.
\item \(1e^{-4}\) (accel bias): \(\sigma = 0.01\) m/s\(^2\). 비교적 큰 불확실 (가속도계 bias 는 잡기 어려움).
\end{itemize}

\textbf{튜닝 가능}: 시퀀스가 시작부터 움직이면 위치/속도/자세 분산을 \(10\times\) 키우는 게 발산 방지에 유리.

\subsection{공개 API: propagate}

\begin{lstlisting}[style=cpp,caption={\texttt{imu\_propagator.cpp:129-149}}]
void ImuPropagator::propagate(double t_new,
                               const Vector3d& gyro_raw,
                               const Vector3d& accel_raw) {
    if (t_prev_ < 0.0) {
        t_prev_ = t_new;
        return;  // 첫 샘플은 dt 모름 → 스킵
    }
    double dt = t_new - t_prev_;
    if (dt <= 0.0 || dt > 0.5) {  // sanity gate
        t_prev_ = t_new;
        return;
    }
    t_prev_ = t_new;

    Vector3d gyro_c  = gyro_raw  - bg;   // bias 보정
    Vector3d accel_c = accel_raw - ba;

    rk4_step(dt, gyro_c, accel_c);          // nominal state 적분
    propagate_cov(dt, gyro_c, accel_c);     // 공분산 전파
}
\end{lstlisting}

\textbf{왜 sanity gate \(dt > 0.5\) 차단?} 시퀀스 시작/종료 경계에서 timestamp 가 튀거나, 데이터 누락이 있을 때 큰 \(dt\) 로 적분하면 \(O(dt^5)\) 라도 거대한 위치 점프가 생긴다. 0.5 s 는 KITTI 의 정상 IMU 간격 0.1 s 의 5배라 안전 마진.

\subsection{공분산 전파의 코드-수식 매핑}

\begin{lstlisting}[style=cpp,caption={\texttt{imu\_propagator.cpp:89-125}}]
void ImuPropagator::propagate_cov(double dt,
                                   const Vector3d& gyro_c,
                                   const Vector3d& accel_c) {
    Mat15 F = Mat15::Zero();
    F.block<3,3>(0,3)   =  Matrix3d::Identity();        // (1)
    F.block<3,3>(3,6)   = -R * so3::hat(accel_c);       // (2)
    F.block<3,3>(3,12)  = -R;                           // (3)
    F.block<3,3>(6,6)   = -so3::hat(gyro_c);            // (4)
    F.block<3,3>(6,9)   = -R;                           // (5)

    Mat15 Phi = Mat15::Identity() + F * dt;             // (6)

    Eigen::Matrix<double, 15, 12> G = ...::Zero();
    G.block<3,3>(3,3)  = R;                             // accel noise -> dv
    G.block<3,3>(6,0)  = R;                             // gyro  noise -> dtheta
    G.block<3,3>(9,6)  = Matrix3d::Identity();          // gyro  walk  -> dbg
    G.block<3,3>(12,9) = Matrix3d::Identity();          // accel walk  -> dba

    double sg=noise_.gyro_noise, sa=noise_.accel_noise;
    double sbg=noise_.gyro_walk, sba=noise_.accel_walk;
    Eigen::Matrix<double, 12, 12> Qc = ...::Zero();
    Qc.block<3,3>(0,0)  = sg*sg  * Matrix3d::Identity();
    Qc.block<3,3>(3,3)  = sa*sa  * Matrix3d::Identity();
    Qc.block<3,3>(6,6)  = sbg*sbg* Matrix3d::Identity();
    Qc.block<3,3>(9,9)  = sba*sba* Matrix3d::Identity();

    Mat15 Qd = (G * Qc * G.transpose()) * dt;           // (7)
    P = Phi * P * Phi.transpose() + Qd;                 // (8)
}
\end{lstlisting}

매핑:
\begin{itemize}
\item (1)\textasciitilde{}(5): 공분산 전파 절의 \(F\) 행렬 5개 비영 블록과 정확히 일치.
\item (6): \(\Phi = I + F dt\) 의 1차 Euler 이산화.
\item (7): \(Q_d = G Q_c G\T dt\). \(\sigma\) 들이 ``밀도'' 단위라 \(dt\) 곱.
\item (8): \(P \leftarrow \Phi P \Phi\T + Q_d\) 의 EKF 표준식.
\end{itemize}

% ---------------------------------------------------------------------
\section{LcEkf — Loosely-Coupled 측정 업데이트}
% ---------------------------------------------------------------------

\subsection{생성자와 lever arm}

\begin{lstlisting}[style=cpp,caption={\texttt{src/ekf/lc\_ekf.cpp:10-18}}]
LcEkf::LcEkf(ImuPropagator& imu,
              const Eigen::Matrix4d& T_cam0_imu,
              double sigma_vo)
    : imu_(imu), sigma_vo_(sigma_vo)
{
    Eigen::Matrix4d T_imu_cam0 = T_cam0_imu.inverse();
    p_IC_ = T_imu_cam0.block<3,1>(0,3);   // body frame 에서 본 cam0 원점
}
\end{lstlisting}

\(T_{\text{cam0,imu}}\) 는 IMU \(\to\) cam0 변환이므로 역행렬이 cam0 \(\to\) IMU. 그 평행이동 부분이 \(\bm{p}_{IC}\) — body frame 에서 cam0 가 어디 있는지 (lever arm).

\subsection{측정 업데이트 전체 코드}

\begin{lstlisting}[style=cpp,caption={\texttt{lc\_ekf.cpp:26-68}}]
bool LcEkf::update_vo(const Vector3d& p_world_cam0) {
    // (1) 예측 측정값
    const Matrix3d& R_wI = imu_.R;
    const Vector3d& p_wI = imu_.p;
    Vector3d h = R_wI * p_IC_ + p_wI;

    // (2) Innovation
    Vector3d innov = p_world_cam0 - h;

    // (3) Measurement Jacobian (3x15)
    Eigen::Matrix<double, 3, 15> H = ...::Zero();
    H.block<3,3>(0,0) = Matrix3d::Identity();
    H.block<3,3>(0,6) = so3::hat(R_wI * p_IC_);

    // (4) Innovation covariance
    Eigen::Matrix3d R_noise = (sigma_vo_*sigma_vo_) * Matrix3d::Identity();
    Eigen::Matrix3d S = H * imu_.P * H.transpose() + R_noise;

    // (5) Kalman gain
    Eigen::Matrix<double, 15, 3> K = imu_.P * H.transpose() * S.inverse();

    // (6) Error correction
    Vec15 dx = K * innov;
    apply_correction(dx);

    // (7) Joseph-form covariance update
    Mat15 I_KH = Mat15::Identity() - K * H;
    imu_.P = I_KH * imu_.P * I_KH.transpose()
           + K * R_noise * K.transpose();

    if (innov.norm() > 5.0)
        std::cerr << "[EKF] Large VO innovation: " << innov.norm() << " m\n";
    return true;
}
\end{lstlisting}

\subsection{상태 보정 적용}

\begin{lstlisting}[style=cpp,caption={\texttt{lc\_ekf.cpp:70-77}}]
void LcEkf::apply_correction(const Vec15& dx) {
    imu_.p  += dx.segment<3>(0);
    imu_.v  += dx.segment<3>(3);
    imu_.R   = imu_.R * so3::Exp(dx.segment<3>(6));
    imu_.R   = Eigen::Quaterniond(imu_.R).normalized().toRotationMatrix();
    imu_.bg += dx.segment<3>(9);
    imu_.ba += dx.segment<3>(12);
}
\end{lstlisting}

\textbf{핵심}: 위치/속도/바이어스는 \(+=\) (가법), 회전만 \(R \cdot \Exp(\delta\bm{\theta})\) 합성. 그 후 정규화로 수치 오차 cosmetics.

% ---------------------------------------------------------------------
\section{StereoTracker — Stereo VO 프론트엔드}
% ---------------------------------------------------------------------

\subsection{생성자: 캘리브레이션과 rectify 맵}

\begin{lstlisting}[style=cpp,caption={\texttt{src/frontend/stereo\_tracker.cpp:16-91}, 핵심부 발췌}]
StereoTracker::StereoTracker(const CameraParams& cam0,
                             const CameraParams& cam1) {
    // (1) Intrinsic 행렬과 distortion 벡터 구성
    K0_ = (cv::Mat_<double>(3,3) << cam0.fx, 0, cam0.cx, ...);
    K1_ = ...

    // (2) cam0 -> cam1 상대 변환
    T_cam1_cam0_ = T_C1_I * T_C0_I.inverse();
    cv::Mat R10(3,3,CV_64F), t10(3,1,CV_64F);  // T_cam1_cam0_ -> R10, t10

    // (3) Stereo rectify (epipolar line 을 수평으로)
    cv::stereoRectify(K0_, dist0_, K1_, dist1_, img_size,
                      R10, t10,
                      R0_rect_, R1_rect_, P0_, P1_, Q,
                      cv::CALIB_ZERO_DISPARITY, -1);

    // (4) 매 픽셀별 undistort+rectify map (런타임에 cv::remap 으로 적용)
    cv::initUndistortRectifyMap(K0_, dist0_, R0_rect_, P0_, img_size,
                                 CV_32FC1, map0x_, map0y_);
    cv::initUndistortRectifyMap(K1_, dist1_, R1_rect_, P1_, img_size,
                                 CV_32FC1, map1x_, map1y_);

    fx_rect_       = P0_.at<double>(0, 0);
    baseline_rect_ = -P1_.at<double>(0, 3) / fx_rect_;
    K_rect_        = P0_(cv::Rect(0, 0, 3, 3)).clone();

    disparity_offset_px_ = -static_cast<float>(fx_rect_ * baseline_rect_ / 3.0);
    orb_     = cv::ORB::create(1500, 1.2f, 8, 15, 0, 2,
                                cv::ORB::HARRIS_SCORE, 31, 7);
    matcher_ = cv::BFMatcher::create(cv::NORM_HAMMING, false);
}
\end{lstlisting}

핵심 결과:
\begin{itemize}
\item \texttt{P0\_}, \texttt{P1\_} : rectified projection matrix. KITTI 의 calib 와 일치해야 함.
\item \texttt{baseline\_rect\_} \(\approx 0.537\) m — 미터 스케일의 출처.
\item \texttt{disparity\_offset\_px\_} = \(-f_x \cdot b / 3\) — guided KLT 의 ``예상 깊이 3m''에 해당하는 시차 초기 추정. 음수인 이유는 cam1 이 우측이라 features 가 \emph{왼쪽으로} 이동.
\end{itemize}

\subsection{한 프레임 처리}

\(N\)번째 프레임에서 \texttt{tracker.process(img\_l, img\_r)} 호출 시:

\begin{enumerate}
\item \textbf{전처리}: BGR \(\to\) gray, \texttt{cv::remap} 으로 stereo rectify 적용.
\item \textbf{Bootstrap (첫 프레임만)}: ORB 로 좌\(\leftrightarrow\)우 매칭, 삼각측량, 깊이 \(0.1 < Z < 50\) m 만 keep, world frame = first cam0.
\item \textbf{시간축 추적}: KLT 로 prev\_left \(\to\) curr\_left.
\item \textbf{스테레오 추적}: guided KLT 로 curr\_left \(\to\) curr\_right (\texttt{disparity\_offset} 으로 초기 추정).
\item \textbf{Outlier 필터}: epipolar 제약 (\(|\Delta v| < 2\) px), disparity 부호 검사.
\item \textbf{삼각측량}: 새 left-right 매칭으로 3D 점 갱신.
\item \textbf{PnP RANSAC}: 보유 3D 점과 curr\_left 픽셀로 카메라 6-DoF pose.
\item \textbf{Gap fill}: 추적 점이 \texttt{target\_features\_=250} 보다 적으면 FAST 검출 + ORB stereo 매칭으로 추가.
\item \textbf{상태 갱신}: prev\_pts \(\leftarrow\) curr\_pts.
\end{enumerate}

% ---------------------------------------------------------------------
\section{KittiRawReader — 데이터 적재}
% ---------------------------------------------------------------------

\subsection{디렉터리 구조 검증과 grayscale 우선}

\begin{lstlisting}[style=cpp,caption={\texttt{src/io/kitti\_raw\_reader.cpp:79-91}}]
std::pair<fs::path, fs::path> select_stereo_roots(const fs::path& root) {
    const std::pair gray{root / "image_00", root / "image_01"};
    if (fs::exists(gray.first / "data") && fs::exists(gray.second / "data"))
        return gray;

    const std::pair color{root / "image_02", root / "image_03"};
    if (fs::exists(color.first / "data") && fs::exists(color.second / "data"))
        return color;

    throw std::runtime_error("Cannot find KITTI stereo image folders ...");
}
\end{lstlisting}

\textbf{왜 grayscale 우선?} 컬러 카메라(\texttt{image\_02/03})는 베이어 보간을 거쳐서 콘트라스트가 흑백 대비 약간 떨어진다. KLT 추적은 강한 그라디언트에 의존하므로 흑백이 안정적. 또한 KITTI 의 흑백 카메라가 더 넓은 다이내믹 레인지.

\subsection{OXTS 파싱: IMU + GT 한 번에}

\begin{lstlisting}[style=cpp,caption={\texttt{kitti\_raw\_reader.cpp:170-205}, 핵심부}]
for (size_t i = 0; i < count; ++i) {
    const auto values = load_numeric_tokens(files[i]);
    if (values.size() < 23) continue;

    const double lat = values[0];
    const double lon = values[1];
    const double alt = values[2];
    const double roll = values[3];
    const double pitch = values[4];
    const double yaw = values[5];

    if (!have_origin) {
        mercator_scale = std::cos(lat * kPi / 180.0);
        origin = Eigen::Vector3d(
            mercator_x(lon, mercator_scale),
            mercator_y(lat, mercator_scale),
            alt);
        have_origin = true;
    }

    ImuData imu;
    imu.timestamp = timestamps[i];
    imu.gyro = Eigen::Vector3d(values[20], values[21], values[22]);  // wf,wl,wu
    imu.accel = Eigen::Vector3d(values[14], values[15], values[16]); // af,al,au
    imu_.push_back(imu);

    GtData gt;
    gt.timestamp = timestamps[i];
    gt.p = Eigen::Vector3d(
        mercator_x(lon, mercator_scale) - origin.x(),
        mercator_y(lat, mercator_scale) - origin.y(),
        alt - origin.z());
    gt.q = Eigen::Quaterniond(rot_from_rpy(roll, pitch, yaw));
    gt.q.normalize();
    gt.v = Eigen::Vector3d(values[7], values[6], values[10]);
    gt_.push_back(gt);
}
\end{lstlisting}

핵심:
\begin{itemize}
\item OXTS .txt 한 줄 30개 값에서 \texttt{values[14..16]} = (af, al, au) 가속도, \texttt{values[20..22]} = (wf, wl, wu) 각속도 \emph{body frame, forward-left-up}.
\item lat/lon/alt 는 첫 패킷을 원점으로 Mercator 투영해 \(\R^3\) 좌표화. (위경도 차이를 미터로 변환)
\item GT 자세는 RPY \(\to\) 쿼터니언.
\item GT 속도는 north/east 가 \texttt{values[6..7]} 이고 east-north-up 순서로 저장 (코드의 \texttt{values[7],[6],[10]} 순서 확인).
\end{itemize}

% ---------------------------------------------------------------------
\section{run\_vio.cpp — 메인 루프}
% ---------------------------------------------------------------------

\subsection{초기화: 정지 IMU 평균}

\begin{lstlisting}[style=cpp,caption={\texttt{apps/run\_vio.cpp:96-155}, 발췌}]
InitialState initialize_from_imu(const std::vector<ImuData>& imu_data,
                                 double t0,
                                 const InitParams& init_params) {
    InitialState init;
    init.t0 = t0;
    init.g_world = Eigen::Vector3d(0.0, 0.0, -init_params.gravity_norm);

    Eigen::Vector3d gyro_sum = Eigen::Vector3d::Zero();
    Eigen::Vector3d accel_sum = Eigen::Vector3d::Zero();
    int sample_count = 0;

    for (const auto& imu : imu_data) {
        if (imu.timestamp < t0) continue;
        if (imu.timestamp > t0 + init_params.static_window_sec) break;
        gyro_sum += imu.gyro;
        accel_sum += imu.accel;
        ++sample_count;
    }
    // ...

    init.bg0 = gyro_sum / sample_count;
    const Vector3d accel_mean = accel_sum / sample_count;
    init.R0 = estimate_initial_orientation(accel_mean);
    return init;
}

Eigen::Matrix3d estimate_initial_orientation(const Vector3d& accel_mean) {
    Eigen::Quaterniond q = Eigen::Quaterniond::FromTwoVectors(
        accel_mean.normalized(),
        Eigen::Vector3d::UnitZ());  // accel_mean -> +Z
    return q.normalized().toRotationMatrix();
}
\end{lstlisting}

\textbf{왜 가속도 평균이 \(+Z\) 방향?} 정지 시 IMU 가속도계는 \(R\T(-\bm{g}) = R\T \cdot (0, 0, +9.81)\T\) 를 출력 (중력의 반대 방향). \(R_0\) 는 이 측정 벡터가 world \(+Z\) 가 되도록 회전을 잡는다.

\textbf{왜 1초 평균?} IMU 잡음은 white 라 평균이 \(\sigma / \sqrt{N}\) 로 줄어든다. KITTI 10 Hz 라 \(N=10\), \(\sqrt{10} \approx 3.2\)배 잡음 감소.

\subsection{메인 처리 루프}

\begin{lstlisting}[style=cpp,caption={\texttt{run\_vio.cpp:310-340}}]
for (const auto& cam : cam_data) {
    if (max_frames > 0 && frame_count >= max_frames) break;

    // (1) IMU 를 cam timestamp 까지 모두 전파
    while (imu_idx < imu_data.size() &&
           imu_data[imu_idx].timestamp <= cam.timestamp) {
        const auto& imu = imu_data[imu_idx++];
        ekf.propagate(imu.timestamp, imu.gyro, imu.accel);
    }

    // (2) 좌/우 png 로드
    const cv::Mat img_l = cv::imread(cam.img_l, cv::IMREAD_COLOR);
    const cv::Mat img_r = cv::imread(cam.img_r, cv::IMREAD_COLOR);
    if (img_l.empty() || img_r.empty()) continue;

    // (3) Stereo VO
    const auto pose = tracker.process(img_l, img_r);

    // (4) VO local -> world 변환
    const Eigen::Vector3d p_world_cam = R_WC0 * pose.t + t_WC0;

    // (5) EKF 측정 보정 (첫 프레임 제외)
    if (pose.valid && frame_count > 0) {
        ekf.update_vo(p_world_cam);
    }

    // (6) 결과 누적
    traj_est.push_back(ekf.position());
    traj_ts.push_back(cam.timestamp);

    Eigen::Vector3d gt_p;
    if (sample_gt_position(gt_data, cam.timestamp, gt_p))
        traj_gt.push_back(gt_p);

    ++frame_count;
}
\end{lstlisting}

\subsection{ATE 계산}

\begin{lstlisting}[style=cpp,caption={\texttt{run\_vio.cpp:182-200}}]
double compute_ate(const std::vector<Vector3d>& est,
                   const std::vector<Vector3d>& gt) {
    if (est.size() != gt.size() || est.empty()) return -1.0;

    Eigen::MatrixXd est_mat(3, est.size());
    Eigen::MatrixXd gt_mat (3, gt.size());
    for (size_t i = 0; i < est.size(); ++i) {
        est_mat.col(i) = est[i];
        gt_mat .col(i) = gt[i];
    }

    const Eigen::Matrix4d T_gt_est =
        Eigen::umeyama(est_mat, gt_mat, /*with_scale=*/false);

    double sum2 = 0.0;
    for (size_t i = 0; i < est.size(); ++i) {
        Eigen::Vector4d p_est(est[i].x(), est[i].y(), est[i].z(), 1.0);
        Eigen::Vector3d aligned = (T_gt_est * p_est).head<3>();
        sum2 += (aligned - gt[i]).squaredNorm();
    }
    return std::sqrt(sum2 / est.size());
}
\end{lstlisting}

\(\texttt{with\_scale=false}\) 로 스테레오의 미터 스케일을 보존하면서 회전+평행이동만 정합.


% =====================================================================
\chapter{운용 절차 — 데이터부터 결과까지}
% =====================================================================

\section{사전 요구사항}

\begin{center}\small
\begin{tabular}{lll}
\toprule
항목 & 권장 사양 & 비고 \\
\midrule
OS & WSL2 Ubuntu 22.04 / Linux & macOS 미검증 \\
컴파일러 & gcc 10+ (C++20) & \texttt{gcc --version} \\
CMake & 3.16+ & \\
디스크 & 단일 drive 0.5\textasciitilde{}5 GB & 전체 raw \(\sim\)180 GB \\
메모리 & 8 GB+ & \\
Python & 3.8+ + numpy/matplotlib & 시각화용 \\
\bottomrule
\end{tabular}
\end{center}

\section{의존성 설치}

\begin{lstlisting}[style=shell]
sudo apt update
sudo apt install -y \
    build-essential cmake git pkg-config \
    libeigen3-dev libopencv-dev libyaml-cpp-dev \
    python3 python3-pip python3-numpy python3-matplotlib
\end{lstlisting}

버전 점검:
\begin{lstlisting}[style=shell]
gcc --version                          # 10.x 이상
pkg-config --modversion eigen3         # 3.3+
pkg-config --modversion opencv4        # 4.x
pkg-config --modversion yaml-cpp       # 0.6+
\end{lstlisting}

\section{KITTI raw 다운로드}

\begin{enumerate}
\item \url{https://www.cvlibs.net/datasets/kitti/raw_data.php} 에서 계정 등록 (이메일 인증 필수).
\item 시퀀스마다 \textbf{2개 zip} 다운:
  \begin{itemize}
  \item \texttt{2011\_09\_26\_calib.zip} (\(\sim\)1 MB) — 날짜당 1회
  \item \texttt{2011\_09\_26\_drive\_XXXX\_sync.zip} — 메인 데이터
  \end{itemize}
\item LiDAR/tracklets/extract 는 받지 말 것 (안 씀).
\end{enumerate}

\subsection{추천 시퀀스}

\begin{center}\small
\begin{tabular}{llll}
\toprule
drive & 카테고리 & 길이 & 권장 용도 \\
\midrule
\texttt{2011\_09\_26\_drive\_0001} & residential & 11 s & 디버그 \\
\texttt{2011\_09\_26\_drive\_0009} & residential & 47 s & 기본 검증 \\
\texttt{2011\_09\_26\_drive\_0117} & city & 1 m 6 s & City 최장 \\
\texttt{2011\_09\_30\_drive\_0028} & residential & 8 m & KITTI Odometry seq 08 \\
\texttt{2011\_10\_03\_drive\_0027} & residential & 7 m & KITTI Odometry seq 00 \\
\bottomrule
\end{tabular}
\end{center}

\subsection{디렉터리 구조}

\begin{lstlisting}[style=shell]
data/kitti_raw/
└── 2011_09_26/
    ├── calib_cam_to_cam.txt
    ├── calib_imu_to_velo.txt
    ├── calib_velo_to_cam.txt
    └── 2011_09_26_drive_0001_sync/
        ├── image_00/{data/, timestamps.txt}
        ├── image_01/, image_02/, image_03/
        └── oxts/{data/, timestamps.txt, dataformat.txt}
\end{lstlisting}

\section{캘리브레이션 추출}

KITTI 의 calib 3종에서 본 프로젝트가 요구하는 \(T_{\text{cam,imu}}\) 4$\times$4 두 개 (cam0, cam1) 를 계산.

\subsection{Chain rule}
\begin{equation}
T_{\text{cam0,imu}} = R_{\text{rect,00}} \cdot T_{\text{cam,velo}} \cdot T_{\text{velo,imu}}
\end{equation}

cam1 은 추가로:
\begin{equation}
T_{\text{cam1,imu}} = T_{\text{cam1,cam0}} \cdot T_{\text{cam0,imu}}
\end{equation}

\(T_{\text{cam1,cam0}}\) 의 평행이동 \(t_x = -\text{baseline}\) (rectified 좌표계에서 cam1 은 cam0 의 \(-X\) 방향).

\subsection{Python 스크립트}

\begin{lstlisting}[style=py,caption={임시 파일로 저장 후 1회 실행}]
from pathlib import Path
import numpy as np

CALIB = Path("/mnt/d/02_research/06_cpp_LC-EKF_VIO_kitti_raw/data/kitti_raw/2011_09_26")

def parse_kv(path):
    out = {}
    for line in path.read_text().splitlines():
        if ":" not in line: continue
        k, v = line.split(":", 1)
        out[k.strip()] = np.fromstring(v, sep=" ")
    return out

def homog(R, t):
    T = np.eye(4); T[:3,:3] = R.reshape(3,3); T[:3,3] = t.reshape(3); return T

cam = parse_kv(CALIB / "calib_cam_to_cam.txt")
v2c = parse_kv(CALIB / "calib_velo_to_cam.txt")
i2v = parse_kv(CALIB / "calib_imu_to_velo.txt")

T_cam_velo = homog(v2c["R"], v2c["T"])
T_velo_imu = homog(i2v["R"], i2v["T"])
R_rect_00  = np.eye(4); R_rect_00[:3,:3] = cam["R_rect_00"].reshape(3,3)

T_cam0_imu = R_rect_00 @ T_cam_velo @ T_velo_imu

P_rect_00 = cam["P_rect_00"].reshape(3,4)
P_rect_01 = cam["P_rect_01"].reshape(3,4)
fx = P_rect_00[0,0]; fy = P_rect_00[1,1]
cx = P_rect_00[0,2]; cy = P_rect_00[1,2]
baseline_x = -P_rect_01[0,3] / fx
T_cam1_cam0 = np.eye(4); T_cam1_cam0[0,3] = -baseline_x
T_cam1_imu = T_cam1_cam0 @ T_cam0_imu

w, h = cam["S_rect_00"].astype(int)
print(f"{int(w)}x{int(h)}, fx={fx:.4f}, fy={fy:.4f}, cx={cx:.4f}, cy={cy:.4f}")
print(f"baseline_x = {baseline_x:.6f}")
\end{lstlisting}

\section{YAML 설정}

\begin{lstlisting}[style=yaml,caption={\texttt{config/kitti\_raw\_2011\_09\_26\_drive\_0001.yaml}}]
dataset_type: "kitti_raw"
dataset: "/mnt/d/.../data/kitti_raw/2011_09_26/2011_09_26_drive_0001_sync"
output:  "/mnt/d/.../results/kitti_raw_2011_09_26_drive_0001"
max_frames: 0           # 0 = 전체

init:
  mode: "stationary_imu"
  static_window_sec: 1.0
  gravity_norm: 9.81

imu_noise:
  gyro:       1.6968e-4
  accel:      2.0000e-3
  gyro_walk:  1.9393e-5
  accel_walk: 3.0000e-3

ekf:
  sigma_vo: 0.10        # VO 측정 잡음 [m]

cam0:
  width: 1242
  height: 375
  fx: 721.5377
  fy: 721.5377
  cx: 609.5593
  cy: 172.8540
  k1: 0.0  # KITTI rect 이미지는 왜곡 0
  k2: 0.0
  p1: 0.0
  p2: 0.0
  T_cam_imu:
    - [ 0.000998747206, -0.999990382071,  0.004259378488, -0.314076869806 ]
    - [ 0.008416901828, -0.004250821136, -0.999955569675,  0.719452035633 ]
    - [ 0.999964048697,  0.001034553284,  0.008412575209, -1.089082940192 ]
    - [ 0.0,             0.0,             0.0,              1.0           ]

cam1:
  # 같은 intrinsics, T_cam_imu 만 baseline 만큼 다름
  ...
\end{lstlisting}

\section{빌드 \& 실행}

\begin{lstlisting}[style=shell]
cd /mnt/d/02_research/06_cpp_LC-EKF_VIO_kitti_raw
mkdir -p build && cd build
cmake .. -DCMAKE_BUILD_TYPE=Release
make -j$(nproc)

cd ..
./build/run_vio config/kitti_raw_2011_09_26_drive_0001.yaml
\end{lstlisting}

정상 stdout (drive\_0001 기준):
\begin{lstlisting}[style=shell]
[KITTI raw] IMU=108  CAM=108  GT=108
Init window samples=11  mean accel= ... 9.8...
Init gravity= 0  0  -9.81
Frame 100  ... tracked=...
========================================
Frames processed : 108
ATE RMSE         : 156.3 cm
========================================
\end{lstlisting}

\section{시각화}

\begin{lstlisting}[style=shell]
python3 tools/plot_trajectory.py results/kitti_raw_2011_09_26_drive_0001/
\end{lstlisting}

생성된 \texttt{trajectory\_plot.png}: 왼쪽 XY top-down (estimated vs GT), 오른쪽 XYZ vs time.


% =====================================================================
\chapter{다중 시퀀스 일괄 테스트}
% =====================================================================

\section{단일 날짜 (sed 복제)}

\begin{lstlisting}[style=shell]
BASE=config/kitti_raw_2011_09_26_drive_0001.yaml
for drv in 0009 0117 0009 0036 0086; do
    NEW=config/kitti_raw_2011_09_26_drive_${drv}.yaml
    sed \
      -e "s|drive_0001_sync|drive_${drv}_sync|g" \
      -e "s|kitti_raw_2011_09_26_drive_0001|kitti_raw_2011_09_26_drive_${drv}|g" \
      "$BASE" > "$NEW"
done
\end{lstlisting}

\section{전체 KITTI raw 자동화 (5개 날짜)}

전체 데이터셋(5 calib 날짜)을 한 번에 다루려면 캘리브레이션 추출 절의 Python 스크립트를 확장해 모든 \texttt{2011\_*\_drive\_*\_sync} 디렉터리에 대해 yaml 자동 생성. (스크립트는 부록 A 참조)

\section{일괄 실행}

\begin{lstlisting}[style=shell]
mkdir -p logs
for cfg in config/kitti_raw_2011_*_drive_*.yaml; do
    name=$(basename "$cfg" .yaml)
    ./build/run_vio "$cfg" 2>&1 | tee "logs/${name}.log"
done
\end{lstlisting}

\section{결과 집계}

\begin{lstlisting}[style=shell]
echo "date,drive,frames,ate_rmse_m" > summary.csv
for m in results/kitti_raw_2011_*_drive_*/metrics.txt; do
    name=$(basename $(dirname "$m"))
    date=$(echo "$name" | sed 's/kitti_raw_//;s/_drive_.*//')
    drv=$(echo  "$name" | sed 's/.*_drive_//')
    f=$(awk -F': ' '/^frames/{print $2}' "$m")
    a=$(awk -F': ' '/^ate_rmse_m/{print $2}' "$m")
    echo "${date},${drv},${f},${a}" >> summary.csv
done
\end{lstlisting}


% =====================================================================
\chapter{결과 해석 \& 트러블슈팅}
% =====================================================================

\section{실패 모드 진단표}

\begin{center}\small
\begin{tabular}{p{0.32\textwidth}p{0.30\textwidth}p{0.32\textwidth}}
\toprule
증상 & 원인 & 처방 \\
\midrule
\texttt{tracked=0} 빈번 & 어두운 시퀀스 / 빠른 회전 / KLT 손실 & grayscale 강제, drive 변경 \\
EKF 발산 (frame 100 후 imu\_p 폭주) & \texttt{sigma\_vo} 너무 작음 & 0.10 \(\to\) 0.20 \\
Init 직후부터 어긋남 & static\_window 안에 차량 출발 & 0.3\textasciitilde{}0.5초로 단축 \\
ATE \texttt{unavailable} & GT 로드 실패 & \texttt{oxts/timestamps.txt} 존재 확인 \\
GT 와 회전된 채 평행이동 & LC-EKF 글로벌 yaw 미관측 & 정상 한계, Tight MSCKF 필요 \\
중간 점프 & VO outlier & \texttt{sigma\_vo} 상향 + chi-square gate 추가 \\
Z축 폭주 & 초기 gravity 추정 실패 & static\_window 단축 \\
모든 프레임 같은 위치 & EKF propagate 누락 & timestamp 단위 확인 \\
\bottomrule
\end{tabular}
\end{center}

\section{LC-EKF 본질적 한계}

\begin{enumerate}
\item VO 누적 드리프트가 EKF 안으로 그대로 들어옴.
\item Yaw 절대 관측 부재 — 글로벌 yaw drift 보정 불가.
\item Feature-level 기하 제약 활용 불가.
\end{enumerate}

이 한계를 넘으려면 Tight MSCKF (OpenVINS 류) 로 가야 한다.


% =====================================================================
\chapter{수정 시나리오 5선}
% =====================================================================

\section{VO 신뢰도 ↑}

\begin{lstlisting}[style=yaml]
# config/...drive_X.yaml
ekf:
  sigma_vo: 0.05    # 0.10 -> 0.05 (잡음 1/2)
\end{lstlisting}

효과: \(R_{\text{noise}}\) 1/4 \(\to\) \(K\) 증가 \(\to\) VO 가 EKF 를 더 강하게 끌어당김. ATE 보통 줄지만 outlier 많으면 발산 위험.

\section{초기 자세 불확실성 ↑}

\begin{lstlisting}[style=cpp]
// src/ekf/imu_propagator.cpp:22
P.block<3,3>(6,6)  = 1e-2 * Matrix3d::Identity();  // 1e-4 -> 1e-2
\end{lstlisting}

효과: 초기 attitude 분산 100배 \(\to\) 첫 EKF 업데이트에서 회전 보정량 큼. 정지 구간에서 흔들림 가능, 1e-3 권장.

\section{특징점 풀 확장}

\begin{lstlisting}[style=cpp]
// include/frontend/stereo_tracker.hpp:57
int target_features_ = 400;  // 250 -> 400
\end{lstlisting}

효과: KLT 추적 풀 확장 \(\to\) PnP 인라이어 두꺼워짐. CPU 시간 비례 증가.

\section{Chi-square outlier gate}

\begin{lstlisting}[style=cpp,caption={\texttt{src/ekf/lc\_ekf.cpp} innovation 직후}]
Vector3d innov = p_world_cam0 - h;

// chi-square 게이트
Eigen::Matrix<double, 3, 15> H_tmp = ...::Zero();
H_tmp.block<3,3>(0,0) = Matrix3d::Identity();
H_tmp.block<3,3>(0,6) = so3::hat(R_wI * p_IC_);
Eigen::Matrix3d S_tmp = H_tmp * imu_.P * H_tmp.transpose()
                     + (sigma_vo_*sigma_vo_) * Matrix3d::Identity();
double chi2 = innov.transpose() * S_tmp.inverse() * innov;
const double GATE_3DOF_99 = 11.345;  // chi^2 (df=3, p=0.99)
if (chi2 > GATE_3DOF_99) {
    std::cerr << "[EKF] gated VO innov: chi2=" << chi2 << "\n";
    return false;
}
\end{lstlisting}

\section{GPS 측정 추가}

\begin{lstlisting}[style=cpp,caption={새 메서드 \texttt{LcEkf::update\_gps}}]
bool LcEkf::update_gps(const Vector3d& p_gps, double sigma_gps) {
    Vector3d h = imu_.p;                        // h(x) = p_wI
    Vector3d innov = p_gps - h;
    Eigen::Matrix<double, 3, 15> H = ...::Zero();
    H.block<3,3>(0,0) = Matrix3d::Identity();   // VO 와 다르게 dtheta 영향 없음
    Eigen::Matrix3d R_n = (sigma_gps*sigma_gps) * Matrix3d::Identity();
    Eigen::Matrix3d S = H * imu_.P * H.transpose() + R_n;
    Eigen::Matrix<double, 15, 3> K = imu_.P * H.transpose() * S.inverse();
    Vec15 dx = K * innov;
    apply_correction(dx);
    Mat15 I_KH = Mat15::Identity() - K * H;
    imu_.P = I_KH * imu_.P * I_KH.transpose() + K * R_n * K.transpose();
    return true;
}
\end{lstlisting}


% =====================================================================
\chapter{Knob 카탈로그}
% =====================================================================

\section{YAML knobs}

\begin{center}\small
\begin{tabular}{lll}
\toprule
knob & 코드 위치 & 영향 \\
\midrule
\texttt{ekf.sigma\_vo} & \texttt{run\_vio.cpp:271} & VO 신뢰도 \\
\texttt{imu\_noise.gyro/accel} & \texttt{run\_vio.cpp:68-74, 269} & 잡음 가정 \\
\texttt{imu\_noise.*\_walk} & 동상 & bias 드리프트 \\
\texttt{init.static\_window\_sec} & \texttt{run\_vio.cpp:80-83, 123-125} & 초기 평균화 \\
\texttt{init.gravity\_norm} & \texttt{run\_vio.cpp:80-83, 117} & 중력 크기 \\
\texttt{max\_frames} & \texttt{run\_vio.cpp:259, 311} & 부분 실행 \\
\texttt{cam0/cam1.fx..p2,T\_cam\_imu} & \texttt{run\_vio.cpp:43-63} & 캘리브레이션 \\
\bottomrule
\end{tabular}
\end{center}

\section{C++ 상수 knobs}

\begin{center}\small
\begin{tabular}{lll}
\toprule
knob & 위치 & 현재값 \\
\midrule
P 초기 분산 (위치) & \texttt{imu\_propagator.cpp:20} & 1e-4 \\
P 초기 분산 (속도) & \texttt{imu\_propagator.cpp:21} & 1e-2 \\
P 초기 분산 (자세) & \texttt{imu\_propagator.cpp:22} & 1e-4 \\
P 초기 분산 (gyro bias) & \texttt{imu\_propagator.cpp:23} & 1e-6 \\
P 초기 분산 (accel bias) & \texttt{imu\_propagator.cpp:24} & 1e-4 \\
dt sanity gate & \texttt{imu\_propagator.cpp:137} & 0.5 s \\
\texttt{target\_features\_} & \texttt{stereo\_tracker.hpp:57} & 250 \\
\texttt{min\_pnp\_inliers\_} & \texttt{stereo\_tracker.hpp:58} & 10 \\
\texttt{disparity\_offset\_px\_} & \texttt{stereo\_tracker.cpp:80} & \(-f_x b/3\) \\
ORB nfeatures & \texttt{stereo\_tracker.cpp:85} & 1500 \\
삼각측량 깊이 게이트 & \texttt{stereo\_tracker.cpp:118-119} & 0.1\textasciitilde{}50 m \\
큰 innov 경고 & \texttt{lc\_ekf.cpp:64-65} & 5.0 m \\
\bottomrule
\end{tabular}
\end{center}


% =====================================================================
\appendix
\chapter{명령어 치트시트}
% =====================================================================

\begin{lstlisting}[style=shell]
# 의존성
sudo apt install -y build-essential cmake libeigen3-dev libopencv-dev libyaml-cpp-dev python3-numpy python3-matplotlib

# 빌드
cd /mnt/d/02_research/06_cpp_LC-EKF_VIO_kitti_raw
mkdir -p build && cd build && cmake .. -DCMAKE_BUILD_TYPE=Release && make -j$(nproc)

# 단일 실행
./run_vio ../config/kitti_raw_2011_09_26_drive_0001.yaml

# 시각화
cd .. && python3 tools/plot_trajectory.py results/kitti_raw_2011_09_26_drive_0001/

# 일괄 실행
mkdir -p logs
for cfg in config/kitti_raw_2011_*_drive_*.yaml; do
    name=$(basename "$cfg" .yaml)
    ./build/run_vio "$cfg" 2>&1 | tee "logs/${name}.log"
done

# 결과 요약
echo "date,drive,frames,ate_rmse_m" > summary.csv
for m in results/kitti_raw_2011_*_drive_*/metrics.txt; do
    name=$(basename $(dirname "$m"))
    d=$(echo "$name" | sed 's/kitti_raw_//;s/_drive_.*//')
    drv=$(echo "$name" | sed 's/.*_drive_//')
    f=$(awk -F': ' '/^frames/{print $2}' "$m")
    a=$(awk -F': ' '/^ate_rmse_m/{print $2}' "$m")
    echo "${d},${drv},${f},${a}" >> summary.csv
done
\end{lstlisting}


% =====================================================================
\chapter{행렬 도감}
% =====================================================================

전체 시스템에 등장하는 모든 주요 행렬:

\begin{center}\small
\begin{longtable}{lllp{0.40\textwidth}}
\toprule
이름 & 크기 & 위치 & 의미 \\
\midrule\endhead
\(T_{\text{cam,imu}}\) & 4$\times$4 & yaml & IMU \(\to\) cam 변환 \\
\(\bm{p}_{IC}\) & 3 & \texttt{lc\_ekf.cpp:17} & body frame 의 cam0 위치 \\
\(R_{WC_0}, \bm{t}_{WC_0}\) & 3$\times$3, 3 & \texttt{run\_vio.cpp:285-288} & VO local \(\to\) world anchor \\
\(R\) & 3$\times$3 & \texttt{imu\_propagator.cpp} & world \(\leftarrow\) IMU body 자세 \\
\(P\) & 15$\times$15 & \texttt{:18-24, :124, lc\_ekf:60} & error-state 공분산 \\
\(F\) & 15$\times$15 & \texttt{imu\_propagator.cpp:93-99} & 연속 전이 행렬 \\
\(G\) & 15$\times$12 & \texttt{:106-110} & 잡음 \(\to\) 상태 매핑 \\
\(Q_c\) & 12$\times$12 & \texttt{:115-119} & 연속 잡음 밀도 \\
\(\Phi = I + Fdt\) & 15$\times$15 & \texttt{:102} & 이산 전이 \\
\(Q_d\) & 15$\times$15 & \texttt{:122} & 이산 잡음 \\
\(H\) & 3$\times$15 & \texttt{lc\_ekf.cpp:41-43} & 측정 야코비안 \\
\(R_{\text{noise}}\) & 3$\times$3 & \texttt{:46} & VO 측정 잡음 \\
\(S\) & 3$\times$3 & \texttt{:47} & innovation 공분산 \\
\(K\) & 15$\times$3 & \texttt{:50} & Kalman gain \\
\(\delta\bm{x}\) & 15 & \texttt{:53} & error-state 보정 \\
\(T_{\text{gt,est}}\) & 4$\times$4 & \texttt{run\_vio.cpp:193} & Umeyama 정합 \\
\(P_0, P_1\) & 3$\times$4 & \texttt{stereo\_tracker.cpp:57-67} & rectified projection \\
\(K_{\text{rect}}\) & 3$\times$3 & 동상 & rectified intrinsics \\
\bottomrule
\end{longtable}
\end{center}


\end{document}
